% !TEX program = lualatex
% ---------------------------------------------------------------------------------------------
%  Readback note 9: The test datum of an auxiliary pair, joint harmonics with staircase weights, and the theta lift as a conjugate descent
%
%  Build:  python3 tools/rb.py note <slug>
%
%  One unit of a Lean artifact, rewritten as ordinary mathematics. The prose is hand-written. Lean names,
%  line ranges and counts come from data.tex; quoted Lean statements are copied from the artifact by the build.
%
%  Rules of the text: standard terminology; every symbol is explained where it first appears and listed again
%  in the notation at the end; no letter has two meanings; a proof is complete or says what it takes for granted;
%  Lean names appear only in the LEAN lines and in the concordance.
%
%  The signature rule (the build enforces it): a function, relation or predicate that the artifact defines is
%  written only through its notation macro, which the build makes from the spec (`notation`) with one slot for
%  each explicit argument of the Lean definition. No argument is fixed and dropped, also not "for brevity" and
%  not for one section. Every argument and every hypothesis of a stated declaration stands inside its
%  statement, or in a Setting that the statement cites by number; only implicit and instance arguments (the
%  ambient objects) may stand in the conventions.
% ---------------------------------------------------------------------------------------------
\documentclass[11pt,a4paper]{article}
\ifdefined\pdfvariable\pdfvariable suppressoptionalinfo \numexpr 1+2+4+8\relax\fi
\usepackage[a4paper,hmargin=35mm,top=22mm,bottom=27mm,footskip=12mm,nohead]{geometry}

\usepackage{fontspec}
\directlua{
  luaotfload.add_fallback("gfallback", {"STIXTwoMath-Regular.otf:mode=node;", "DejaVuSans.ttf:mode=node;"})
  luaotfload.add_fallback("gmonofallback", {"DejaVuSansMono.ttf:mode=node;", "STIXTwoMath-Regular.otf:mode=node;"})
}
\setmainfont{NewCM10}[Extension=.otf,UprightFont=*-Book,ItalicFont=*-BookItalic,
  BoldFont=*-Bold,BoldItalicFont=*-BoldItalic,SmallCapsFont=*-Book,SmallCapsFeatures={Letters=SmallCaps},
  RawFeature={fallback=gfallback}]
\setsansfont{NewCMSans10}[Extension=.otf,UprightFont=*-Book,ItalicFont=*-BookOblique,
  BoldFont=*-Bold,BoldItalicFont=*-BoldOblique,RawFeature={fallback=gfallback}]
\setmonofont{JetBrainsMono}[Extension=.otf,UprightFont=*-Regular,ItalicFont=*-Italic,
  BoldFont=*-SemiBold,BoldItalicFont=*-SemiBoldItalic,Scale=MatchLowercase,
  RawFeature={-calt;-liga;fallback=gmonofallback}]
\usepackage{amsmath}
\usepackage{amsthm}
\usepackage[math-style=TeX]{unicode-math}
\setmathfont{NewCMMath-Book.otf}
\emergencystretch=2.5em\relax

\usepackage{xcolor}
\usepackage{tikz}
\usetikzlibrary{calc,arrows.meta,positioning,backgrounds}
\usepackage[protrusion=true,expansion=true]{microtype}
\usepackage{enumitem}
\usepackage{needspace}
\usepackage{booktabs}
\usepackage{array}
\usepackage{longtable}
\usepackage{fancyhdr}
\usepackage{lastpage}
\usepackage{multicol}
\usepackage[hidelinks,pdfusetitle]{hyperref}

\definecolor{ink}{HTML}{0F172A}
\definecolor{inksoft}{HTML}{475569}
\definecolor{inkfaint}{HTML}{94A3B8}
\definecolor{rulegrey}{HTML}{CBD5E1}
\definecolor{arith}{HTML}{C2410C}
\definecolor{est}{HTML}{7C3AED}
\definecolor{ana}{HTML}{0891B2}
\definecolor{ok}{HTML}{047857}
\definecolor{leancol}{HTML}{3730A3}

\title{The test datum of an auxiliary pair, joint harmonics with staircase weights, and the theta lift as a conjugate descent}
\author{grh-compact}
\date{}

% ------------------------------------------------------------------------------------- macros
\newcommand{\leanfont}{\ttfamily\fontsize{8.6}{10.5}\selectfont\color{leancol}}
\newcommand{\lean}[1]{\csname LN@#1\endcsname}
\newcommand{\fillin}[1]{{\color{arith}\itshape[#1]}}  % a place of the template that is not written yet; the build reports every one that is left
\newcommand{\leanquote}[1]{\csname LQ@#1\endcsname}
% generated by rb.py note; do not edit
\newcommand{\ConcDecls}{15}
\newcommand{\Pin}{72fead10}
\newcommand{\PinText}{that of commit {\leanfont 72fead10}}
\newcommand{\ArtifactName}{grh-compact}
\expandafter\newcommand\csname LN@TwistGL\endcsname{\href{https://github.com/lsernau/grh-lean/blob/72fead102f663d5c7a407964a13dc0cf9d977207/grh-compact/src/New/Forms/AuxiliaryPair.lean\#L15-L16}{{\leanfont LevelPlaces.\allowbreak{}TwistGL}}}
\expandafter\newcommand\csname LN@twistLift\endcsname{\href{https://github.com/lsernau/grh-lean/blob/72fead102f663d5c7a407964a13dc0cf9d977207/grh-compact/src/New/Forms/AuxiliaryPair.lean\#L24-L25}{{\leanfont LevelPlaces.\allowbreak{}twistLift}}}
\expandafter\newcommand\csname LN@levelIndicator\endcsname{\href{https://github.com/lsernau/grh-lean/blob/72fead102f663d5c7a407964a13dc0cf9d977207/grh-compact/src/New/Forms/AuxiliaryPair.lean\#L29-L32}{{\leanfont AuxPair.\allowbreak{}levelIndicator}}}
\expandafter\newcommand\csname LN@upsiloneqsum\endcsname{\href{https://github.com/lsernau/grh-lean/blob/72fead102f663d5c7a407964a13dc0cf9d977207/grh-compact/src/New/Forms/AuxiliaryPair.lean\#L39-L50}{{\leanfont AuxPair.\allowbreak{}upsilon\_eq\_sum}}}
\expandafter\newcommand\csname LN@pairWeilupsilon\endcsname{\href{https://github.com/lsernau/grh-lean/blob/72fead102f663d5c7a407964a13dc0cf9d977207/grh-compact/src/New/Forms/AuxiliaryPair.lean\#L55-L67}{{\leanfont AuxPair.\allowbreak{}pairWeil\_upsilon}}}
\expandafter\newcommand\csname LN@paramV\endcsname{\href{https://github.com/lsernau/grh-lean/blob/72fead102f663d5c7a407964a13dc0cf9d977207/grh-compact/src/New/Forms/AuxiliaryPair.lean\#L74-L74}{{\leanfont AuxPair.\allowbreak{}paramV}}}
\expandafter\newcommand\csname LN@isHCParamparamV\endcsname{\href{https://github.com/lsernau/grh-lean/blob/72fead102f663d5c7a407964a13dc0cf9d977207/grh-compact/src/New/Forms/AuxiliaryPair.lean\#L79-L79}{{\leanfont AuxPair.\allowbreak{}isHCParam\_paramV}}}
\expandafter\newcommand\csname LN@FitsParams\endcsname{\href{https://github.com/lsernau/grh-lean/blob/72fead102f663d5c7a407964a13dc0cf9d977207/grh-compact/src/New/Forms/AuxiliaryPair.lean\#L85-L90}{{\leanfont AuxPair.\allowbreak{}FitsParams}}}
\expandafter\newcommand\csname LN@realPairWeilfockLetters\endcsname{\href{https://github.com/lsernau/grh-lean/blob/72fead102f663d5c7a407964a13dc0cf9d977207/grh-compact/src/New/Forms/AuxiliaryPair.lean\#L97-L101}{{\leanfont realPairWeil\_fockLetters}}}
\expandafter\newcommand\csname LN@realPairWeilharmonicPoly\endcsname{\href{https://github.com/lsernau/grh-lean/blob/72fead102f663d5c7a407964a13dc0cf9d977207/grh-compact/src/New/Forms/AuxiliaryPair.lean\#L106-L115}{{\leanfont AuxPair.\allowbreak{}realPairWeil\_harmonicPoly}}}
\expandafter\newcommand\csname LN@existsisJointHarmonic\endcsname{\href{https://github.com/lsernau/grh-lean/blob/72fead102f663d5c7a407964a13dc0cf9d977207/grh-compact/src/New/Forms/AuxiliaryPair.lean\#L120-L138}{{\leanfont AuxPair.\allowbreak{}FitsParams.\allowbreak{}exists\_isJointHarmonic}}}
\expandafter\newcommand\csname LQ@existsisJointHarmonic\endcsname{\leanquoteline{0}{lemma AuxPair.FitsParams.exists\_isJointHarmonic \{l\} (h : O.FitsParams l) : ∃ Φ : \_ → \_,}\leanquoteline{4}{(∀ w, IsJointHarmonic (.std 1 O.signs (negCount O.signs) (-((card n / 2 : ℕ) + O.weight w)))}\leanquoteline{2}{(O.paramV w) (-l w) (Φ w)) ∧ O.harmonicPoly (highestVector \_ O.paramV) = ∏ w,}\leanquoteline{2}{(Φ w).rename (w, ·) ∧ ∀ w c, (toHCParam O.signs (-l w)).lowestKType c = staircaseWeight n O.signs}\leanquoteline{2}{O.ord (if · = 1 then O.colsPlus w else O.colsMinus w) (-((card n / 2 : ℕ) + O.weight w)) c}}
\expandafter\newcommand\csname LN@isEdge\endcsname{\href{https://github.com/lsernau/grh-lean/blob/72fead102f663d5c7a407964a13dc0cf9d977207/grh-compact/src/New/Forms/AuxiliaryPair.lean\#L145-L148}{{\leanfont AuxPair.\allowbreak{}FitsParams.\allowbreak{}isEdge}}}
\expandafter\newcommand\csname LN@pairWeilfockDatum\endcsname{\href{https://github.com/lsernau/grh-lean/blob/72fead102f663d5c7a407964a13dc0cf9d977207/grh-compact/src/New/Forms/AuxiliaryPair.lean\#L153-L155}{{\leanfont AuxPair.\allowbreak{}pairWeil\_fockDatum}}}
\expandafter\newcommand\csname LN@conjKernelequivariant\endcsname{\href{https://github.com/lsernau/grh-lean/blob/72fead102f663d5c7a407964a13dc0cf9d977207/grh-compact/src/New/Forms/AuxiliaryPair.lean\#L160-L164}{{\leanfont AuxPair.\allowbreak{}conjKernel\_equivariant}}}
\expandafter\newcommand\csname LN@thetaeqsmulstarliftswap\endcsname{\href{https://github.com/lsernau/grh-lean/blob/72fead102f663d5c7a407964a13dc0cf9d977207/grh-compact/src/New/Forms/AuxiliaryPair.lean\#L169-L179}{{\leanfont AuxPair.\allowbreak{}theta\_eq\_smul\_star\_lift\_swap}}}
\expandafter\newcommand\csname LQ@thetaeqsmulstarliftswap\endcsname{\leanquoteline{0}{lemma AuxPair.theta\_eq\_smul\_star\_lift\_swap Ξ : O.θ Ξ = (card (WeylIndex n O.κ) * stdLevelIndex}\leanquoteline{2}{O.levelGroup : ℂ) • star (Space.lift O.pairThetaKernel.swap}\leanquoteline{2}{(O.fockDatum (highestVector \_ O.paramV)) (star (ofAlgForm O.paramV 1 Ξ.1)))}}
\newcommand{\ConcordanceRows}{\crow{TwistGL}{\href{https://github.com/lsernau/grh-lean/blob/72fead102f663d5c7a407964a13dc0cf9d977207/grh-compact/src/New/Forms/AuxiliaryPair.lean\#L15-L16}{{\leanfont LevelPlaces.\allowbreak{}TwistGL}}}{AuxiliaryPair.lean}{15--16}
\crow{twistLift}{\href{https://github.com/lsernau/grh-lean/blob/72fead102f663d5c7a407964a13dc0cf9d977207/grh-compact/src/New/Forms/AuxiliaryPair.lean\#L24-L25}{{\leanfont LevelPlaces.\allowbreak{}twistLift}}}{AuxiliaryPair.lean}{24--25}
\crow{levelIndicator}{\href{https://github.com/lsernau/grh-lean/blob/72fead102f663d5c7a407964a13dc0cf9d977207/grh-compact/src/New/Forms/AuxiliaryPair.lean\#L29-L32}{{\leanfont AuxPair.\allowbreak{}levelIndicator}}}{AuxiliaryPair.lean}{29--32}
\crow{upsiloneqsum}{\href{https://github.com/lsernau/grh-lean/blob/72fead102f663d5c7a407964a13dc0cf9d977207/grh-compact/src/New/Forms/AuxiliaryPair.lean\#L39-L50}{{\leanfont AuxPair.\allowbreak{}upsilon\_eq\_sum}}}{AuxiliaryPair.lean}{39--50}
\crow{pairWeilupsilon}{\href{https://github.com/lsernau/grh-lean/blob/72fead102f663d5c7a407964a13dc0cf9d977207/grh-compact/src/New/Forms/AuxiliaryPair.lean\#L55-L67}{{\leanfont AuxPair.\allowbreak{}pairWeil\_upsilon}}}{AuxiliaryPair.lean}{55--67}
\crow{paramV}{\href{https://github.com/lsernau/grh-lean/blob/72fead102f663d5c7a407964a13dc0cf9d977207/grh-compact/src/New/Forms/AuxiliaryPair.lean\#L74-L74}{{\leanfont AuxPair.\allowbreak{}paramV}}}{AuxiliaryPair.lean}{74--74}
\crow{isHCParamparamV}{\href{https://github.com/lsernau/grh-lean/blob/72fead102f663d5c7a407964a13dc0cf9d977207/grh-compact/src/New/Forms/AuxiliaryPair.lean\#L79-L79}{{\leanfont AuxPair.\allowbreak{}isHCParam\_paramV}}}{AuxiliaryPair.lean}{79--79}
\crow{FitsParams}{\href{https://github.com/lsernau/grh-lean/blob/72fead102f663d5c7a407964a13dc0cf9d977207/grh-compact/src/New/Forms/AuxiliaryPair.lean\#L85-L90}{{\leanfont AuxPair.\allowbreak{}FitsParams}}}{AuxiliaryPair.lean}{85--90}
\crow{realPairWeilfockLetters}{\href{https://github.com/lsernau/grh-lean/blob/72fead102f663d5c7a407964a13dc0cf9d977207/grh-compact/src/New/Forms/AuxiliaryPair.lean\#L97-L101}{{\leanfont realPairWeil\_fockLetters}}}{AuxiliaryPair.lean}{97--101}
\crow{realPairWeilharmonicPoly}{\href{https://github.com/lsernau/grh-lean/blob/72fead102f663d5c7a407964a13dc0cf9d977207/grh-compact/src/New/Forms/AuxiliaryPair.lean\#L106-L115}{{\leanfont AuxPair.\allowbreak{}realPairWeil\_harmonicPoly}}}{AuxiliaryPair.lean}{106--115}
\crow{existsisJointHarmonic}{\href{https://github.com/lsernau/grh-lean/blob/72fead102f663d5c7a407964a13dc0cf9d977207/grh-compact/src/New/Forms/AuxiliaryPair.lean\#L120-L138}{{\leanfont AuxPair.\allowbreak{}FitsParams.\allowbreak{}exists\_isJointHarmonic}}}{AuxiliaryPair.lean}{120--138}
\crow{isEdge}{\href{https://github.com/lsernau/grh-lean/blob/72fead102f663d5c7a407964a13dc0cf9d977207/grh-compact/src/New/Forms/AuxiliaryPair.lean\#L145-L148}{{\leanfont AuxPair.\allowbreak{}FitsParams.\allowbreak{}isEdge}}}{AuxiliaryPair.lean}{145--148}
\crow{pairWeilfockDatum}{\href{https://github.com/lsernau/grh-lean/blob/72fead102f663d5c7a407964a13dc0cf9d977207/grh-compact/src/New/Forms/AuxiliaryPair.lean\#L153-L155}{{\leanfont AuxPair.\allowbreak{}pairWeil\_fockDatum}}}{AuxiliaryPair.lean}{153--155}
\crow{conjKernelequivariant}{\href{https://github.com/lsernau/grh-lean/blob/72fead102f663d5c7a407964a13dc0cf9d977207/grh-compact/src/New/Forms/AuxiliaryPair.lean\#L160-L164}{{\leanfont AuxPair.\allowbreak{}conjKernel\_equivariant}}}{AuxiliaryPair.lean}{160--164}
\crow{thetaeqsmulstarliftswap}{\href{https://github.com/lsernau/grh-lean/blob/72fead102f663d5c7a407964a13dc0cf9d977207/grh-compact/src/New/Forms/AuxiliaryPair.lean\#L169-L179}{{\leanfont AuxPair.\allowbreak{}theta\_eq\_smul\_star\_lift\_swap}}}{AuxiliaryPair.lean}{169--179}}
% the notations of the defined functions: one slot for each explicit argument of the Lean declaration
\newcommand{\TwistGL}[2]{\mathrm{GL}^{\mathrm{tw}}_{#2}(#1)}
\newcommand{\twistLift}[2]{\operatorname{lift}_{#1}(#2)}
\newcommand{\levelIndicator}[2]{\mathbf{1}_{#1}(#2)}
\newcommand{\paramV}[3]{l^{V}_{#1}(#2,#3)}
\newcommand{\FitsParams}[2]{#2\ \text{fits}\ #1}
\newcommand{\UnitMembers}{16}
\newcommand{\UnitInNote}{15}
\newcommand{\UnitProofs}{7}
\newcommand{\UnitSteps}{29}
\newcommand{\UnitAllSteps}{31}


\linespread{1.08}                       % 11 pt on 14.7 pt
\setlength{\parindent}{0pt}
\setlength{\parskip}{5.5pt plus 1.5pt minus 0.5pt}
\setlength{\abovedisplayskip}{7.5pt plus 2pt minus 3pt}
\setlength{\belowdisplayskip}{7.5pt plus 2pt minus 3pt}
\setlength{\abovedisplayshortskip}{3pt plus 2pt}
\setlength{\belowdisplayshortskip}{6pt plus 2pt minus 2pt}
\emergencystretch=1.5em
\clubpenalty=10000 \widowpenalty=10000 \displaywidowpenalty=10000
\relpenalty=7000 \binoppenalty=9000
\raggedbottom
\pagestyle{fancy}\fancyhf{}\renewcommand{\headrulewidth}{0pt}
\fancyfoot[C]{\sffamily\fontsize{7.6}{9}\selectfont\color{inkfaint}\thepage\,/\,\pageref*{LastPage}}
\setlist{topsep=2pt,partopsep=0pt,parsep=0pt,itemsep=4pt,leftmargin=2.1em,labelsep=0.55em}

\newcommand{\Q}{\mathbb{Q}}
\newcommand{\Z}{\mathbb{Z}}
\newcommand{\F}{\mathbb{F}}
\newcommand{\len}{\operatorname{len}}
\newcommand{\cok}{\operatorname{cok}}
\newcommand{\End}{\operatorname{End}}
\newcommand{\Hom}{\operatorname{Hom}}
\newcommand{\Gal}{\operatorname{Gal}}
\newcommand{\diag}{\operatorname{diag}}
\newcommand{\kw}[1]{\emph{#1}}
\newcommand{\m}[1]{\mbox{$#1$}}        % a short formula that must not be broken across lines
\newcommand{\head}[1]{\par\addvspace{17pt plus 4pt minus 2pt}\Needspace{6\baselineskip}%
  \noindent{\fontsize{13.2}{16}\selectfont\bfseries\color{ink}#1}\par\nopagebreak\vspace{-1pt}}
\newcommand{\sect}[2]{\refstepcounter{section}\label{#2}\par\addvspace{19pt plus 4pt minus 2pt}\Needspace{7\baselineskip}%
  \noindent\llap{{\fontsize{13.2}{16}\selectfont\bfseries\color{inkfaint}\thesection}\hspace{9pt}}%
  {\fontsize{13.2}{16}\selectfont\bfseries\color{ink}#1}\par\nopagebreak\vspace{-1pt}}
\newcommand{\subhead}[1]{\par\addvspace{9pt plus 2pt}\Needspace{4\baselineskip}\noindent\textbf{#1}\enspace\ignorespaces}
% where a statement lives in the Lean archive
\newcommand{\leantag}{{\sffamily\bfseries\fontsize{6.3}{8}\selectfont\color{leancol}\addfontfeatures{LetterSpace=7}LEAN}}
\newcommand{\sep}{\unskip\nobreak\hspace{5pt}{\color{inkfaint}\textperiodcentered}\hspace{5pt plus 2pt}\ignorespaces}
\newcommand{\leanline}[1]{\par\nopagebreak\vspace{-2.5pt}{\fontsize{8.9}{12.4}\selectfont\color{inksoft}\raggedright
  \noindent\llap{\leantag\hspace{9pt}}#1\par}\vspace{1pt}}
% a displayed statement in a tinted box: #1 colour, #2 text
\newcommand{\statement}[2]{\par\addvspace{7pt}\begingroup\setlength{\fboxsep}{8.5pt}\setlength{\fboxrule}{0.6pt}%
  \noindent\fcolorbox{#1!30!white}{#1!3!white}{\parbox{\dimexpr\linewidth-2\fboxsep-2\fboxrule\relax}{#2}}%
  \endgroup\par\addvspace{7pt}}
\newcommand{\nt}[2]{\par\noindent\hangindent=1.1em\hangafter=1 #1\enspace{\color{inksoft}#2}\par\vspace{1pt}}
\newcommand{\legdot}[1]{\tikz[baseline=-0.62ex]\fill[#1] (0,0) circle (2.5pt);\hspace{3pt}}
\newcommand{\crow}[4]{\csname CS@#1\endcsname & #2 & \textcolor{inksoft}{#3} & \textcolor{inksoft}{#4}\\}
\newcommand{\leanquoteline}[2]{\par\noindent\hangindent=\dimexpr#1\fontdimen2\font+2.2em\relax\hangafter=1
  \hspace*{#1\fontdimen2\font}#2}
\newcommand{\quoted}[1]{\par\vspace{3pt}{\ttfamily\fontsize{7.5}{10.4}\selectfont\color{leancol}\raggedright
  \setlength{\parskip}{0pt}\leanquote{#1}\par}\vspace{4pt}}
\newcommand{\bx}[3]{{\sffamily\bfseries\fontsize{7.3}{9.2}\selectfont\color{#1!82!black}#2}\\[1.5pt]#3}
\newcommand{\cs}[2]{\expandafter\newcommand\csname CS@#1\endcsname{#2}}

% theorems
\newtheoremstyle{plainx}{7pt plus 2pt}{7pt plus 2pt}{\itshape}{}{\bfseries}{.}{0.5em}{}
\newtheoremstyle{uprightx}{7pt plus 2pt}{7pt plus 2pt}{}{}{\bfseries}{.}{0.5em}{}
\theoremstyle{plainx}
\newtheorem{theorem}{Theorem}[section]
\newtheorem{lemma}[theorem]{Lemma}
\newtheorem{proposition}[theorem]{Proposition}
\newtheorem{corollary}[theorem]{Corollary}
\theoremstyle{uprightx}
\newtheorem{definition}[theorem]{Definition}
\newtheorem{setting}[theorem]{Setting}          % hypotheses shared by several statements; each of them cites it by number
\newtheorem{remark}[theorem]{Remark}
\newtheorem{example}[theorem]{Example}
\renewcommand{\qedsymbol}{{\setlength{\fboxsep}{0pt}\setlength{\fboxrule}{0.5pt}\fbox{\rule{0pt}{0.46em}\hspace{0.46em}}}}
% what holds throughout the note: the ambient objects only (implicit and instance arguments), never a hypothesis
\newenvironment{conventions}{\par\addvspace{5pt}\noindent\textbf{Conventions.}\enspace\ignorespaces}{\par\addvspace{3pt}}
\makeatletter
\renewenvironment{proof}[1][Proof]{\par\pushQED{\qed}\normalfont\topsep0pt\partopsep0pt
  \trivlist\item[\hskip\labelsep\itshape #1.]\ignorespaces}{\popQED\endtrivlist\@endpefalse}
\makeatother

% where each declaration of the concordance is treated: \cs{key}{place}, the place naming a \label or a lettered item
\cs{TwistGL}{Definition~\ref{d:twist}}
\cs{twistLift}{Definition~\ref{d:twist}}
\cs{levelIndicator}{Definition~\ref{d:indicator}}
\cs{upsiloneqsum}{Proposition~\ref{p:sum}}
\cs{pairWeilupsilon}{Proposition~\ref{p:law}}
\cs{paramV}{Definition~\ref{d:paramV}}
\cs{isHCParamparamV}{Lemma~\ref{l:hc}}
\cs{FitsParams}{Definition~\ref{d:fits}}
\cs{realPairWeilfockLetters}{Lemma~\ref{l:letters}}
\cs{realPairWeilharmonicPoly}{Lemma~\ref{l:har}}
\cs{existsisJointHarmonic}{Theorem 2}
\cs{isEdge}{Corollary~\ref{c:edge}}
\cs{pairWeilfockDatum}{Lemma~\ref{l:datum}}
\cs{conjKernelequivariant}{Lemma~\ref{l:kernel}}
\cs{thetaeqsmulstarliftswap}{Theorem 1}
% notation of this note; objects of other units and of the background library are written with all their arguments
\newcommand{\CC}{\mathbb{C}}
\newcommand{\ZZ}{\mathbb{Z}}
\newcommand{\AAf}{\mathbb{A}_{\mathrm{f}}}
\newcommand{\FF}{\mathbb{F}}
\newcommand{\oo}{\mathfrak{o}}
\newcommand{\ohat}{\hat{\mathfrak{o}}}
\newcommand{\Ginf}[1]{G_{\infty}(#1)}
\newcommand{\Gfin}[1]{G_{\mathrm{f}}(#1)}
\newcommand{\Gad}[1]{G(#1)}
\newcommand{\Kz}[1]{\mathbf{K}_{0}(#1)}
\newcommand{\std}[2]{\mathrm{I}_{#2}(#1)}
\newcommand{\one}[1]{1_{#1}}
\newcommand{\WO}[1]{W(#1)}
\newcommand{\sgO}[1]{\varepsilon_{#1}}
\newcommand{\rmi}[1]{r_{-}(#1)}
\newcommand{\rpl}[1]{r_{+}(#1)}
\newcommand{\incO}[1]{\operatorname{inc}_{#1}}
\newcommand{\IncM}[1]{\operatorname{Inc}_{#1}}
\newcommand{\ordO}[2]{\operatorname{ord}_{#1}(#2)}
\newcommand{\kp}[3]{\kappa^{+}_{#1}(#2,#3)}
\newcommand{\km}[3]{\kappa^{-}_{#1}(#2,#3)}
\newcommand{\kap}[1]{\kappa_{#1}}
\newcommand{\eO}[1]{e_{#1}}
\newcommand{\muO}[1]{\mu_{#1}}
\newcommand{\wtO}[2]{\operatorname{wt}_{#1}(#2)}
\newcommand{\psiO}[2]{\psi_{#1,#2}}
\newcommand{\MO}[1]{M_{#1}}
\newcommand{\QO}[1]{Q_{#1}}
\newcommand{\SigO}[1]{\Sigma_{#1}}
\newcommand{\ChGL}[2]{\mathrm{GL}^{\mathrm{ch}}_{#2}(#1)}
\newcommand{\chO}[1]{\chi_{#1}}
\newcommand{\chpO}[1]{\chi'_{#1}}
\newcommand{\omO}[1]{\omega_{#1}}
\newcommand{\thO}[1]{\theta_{#1}}
\newcommand{\SO}[1]{\mathcal{S}(#1)}
\newcommand{\LO}[1]{\mathbf{L}(#1)}
\newcommand{\KO}[1]{\mathbf{K}(#1)}
\newcommand{\etaO}[1]{\eta_{#1}}
\newcommand{\tetaO}[1]{\tilde{\eta}_{#1}}
\newcommand{\Ups}[1]{\Upsilon(#1)}
\newcommand{\PsiO}[2]{\Psi_{#1}(#2)}
\newcommand{\Har}[2]{\operatorname{Har}_{#1}(#2)}
\newcommand{\JO}[2]{\operatorname{J}_{#1}(#2)}
\newcommand{\cK}[3]{\mathcal{C}_{#1}(#2,#3)}
\newcommand{\Th}[2]{\operatorname{Th}_{#1}(#2)}
\newcommand{\FormO}[1]{\operatorname{Form}(#1)}
\newcommand{\Vk}[2]{\mathbb{V}_{#1,#2}}
\newcommand{\rhoW}[2]{\rho_{#1,#2}}
\newcommand{\vplus}[2]{v^{+}_{#1}(#2)}
\newcommand{\gram}[3]{\langle #2,#3\rangle_{#1}}
\newcommand{\algF}[3]{F_{#1,#2}(#3)}
\newcommand{\Desc}[4]{\operatorname{Desc}^{#1,#2}_{#3,#4}}
\newcommand{\RPW}[2]{\Omega_{#1,#2}}
\newcommand{\pvO}[2]{v_{#1,#2}}
\newcommand{\FockO}[1]{\mathcal{F}(#1)}
\newcommand{\FockP}[2]{\mathcal{F}(#1,#2)}
\newcommand{\pvP}[3]{v^{#1,#2}_{#3}}
\newcommand{\Let}[3]{\operatorname{Let}_{#1,#2,#3}}
\newcommand{\Pol}[1]{\operatorname{Pol}_{#1}}
\newcommand{\Sch}[1]{S_{#1}}
\newcommand{\tk}[2]{\mathbf{e}_{#1,#2}}
\newcommand{\lkt}[2]{\Lambda_{#1}(#2)}
\newcommand{\stw}[6]{\operatorname{sw}_{#1}(#2,#3,#4,#5)(#6)}
\newcommand{\modD}[5]{#5\otimes[#4]^{#1,#2}_{#3}}
\newcommand{\lmod}[2]{\mathfrak{m}_{#2}(#1)}
\newcommand{\red}[1]{\operatorname{red}_{#1}}
\newcommand{\heisf}[1]{\varrho(#1)}
\newcommand{\lnm}[1]{{\leanfont #1}}
\newcommand{\conj}[1]{\overline{#1}}
\newcommand{\dd}{\mathrm{d}}

\begin{document}
\color{ink}

{\raggedright
{\sffamily\bfseries\fontsize{7.2}{9}\selectfont\color{leancol}\addfontfeatures{LetterSpace=8}READBACK NOTE 9\par}
\vspace{5pt}
{\fontsize{21.5}{25.5}\selectfont\bfseries The test datum of an auxiliary pair, joint harmonics with staircase weights, and the theta lift as a conjugate descent\par}
\vspace{5pt}
{\fontsize{12.4}{16}\selectfont\itshape\color{inksoft} One unit of the Lean artifact \ArtifactName, read back into
mathematics\par}}
\vspace{8pt}
{\color{rulegrey}\hrule height 0.6pt}
\vspace{7pt}
\thispagestyle{empty}

\textbf{Summary.} An \emph{auxiliary pair} is a dual pair of unitary groups over a CM field, an indefinite space~$W$
of rank~$r$ and the standard definite space of rank \m{n>r}, with a level at two families of split places and
a Hecke character; it carries one test datum for the Weil representation, the tensor product of the Fock vector of a
harmonic polynomial and a finite datum~$\Upsilon$, and a map from algebraic forms on the definite group to
functions on the indefinite one, defined as a finite sum over classes. This unit proves the three facts about
this datum that the rest of Theorem~C's proof uses. The finite datum is a sum of translates of an indicator
function and transforms under the level by an explicit character. The polynomial at the highest vector is a
product, over the infinite places, of products of minors (staircases), which are joint harmonics for matched
discrete series parameters. And the class sum is a fixed multiple of the complex conjugate of a theta descent.

\statement{leancol}{\fontsize{9.6}{13.2}\selectfont\textbf{What this note is.} The artifact \ArtifactName{} is a
body of Lean code: the compact edition of a formal proof, due to L.~Alpöge with the AI system Claude, whose final
theorems are the Riemann hypothesis and its generalisation to Dirichlet \emph{L}-functions. That headline is a
claim. The authors report that every file compiles and that Lean's kernel accepted the proofs; an independent
build is recorded in the handout's record of checks when it reports. No mathematician has read the argument
through, and nothing here says that the Riemann hypothesis is proved. This note treats one unit of the artifact:
\UnitInNote{} of its \UnitMembers{} declarations, all that have a name (the remaining one is the remark that a
finite product of finite groups is finite). The proofs are given in full, relative to two kinds of inputs, which
are stated and not proved: the definition of an auxiliary pair and of its objects, which belongs to the unit on
the rows (Setting~\ref{s:aux}); and results of the background library on the Weil representation, on joint
harmonics and on algebraic forms (Setting~\ref{s:weil} and the facts (Q1) to (Q11)). Every statement was compared
with its Lean text, argument by argument. The Lean text read is \PinText; Appendix~A is the concordance of
\ConcDecls{} declarations. This is the second of three notes on the theta correspondence behind Theorem~C; it
uses the Fock space of the note \emph{Decay of Heisenberg matrix coefficients of homogeneous Fock vectors,
uniform in the degree}.}

\sect{The statements}{s:intro}
The objects are defined in §\ref{s:setting}. In brief: for an auxiliary pair~$O$, $\WO{O}$ is its indefinite
space, with signs $\sgO{O}$ on the index set~$[r]$, and $\std{E}{n}$ the definite one; $\FormO{O}$ is a space of
algebraic modular forms on the definite unitary group, with values in a representation
\m{\rhoW{\kap{O}}{\eO{O}}} on a space \m{\Vk{n}{\kap{O}}}; \m{\Th{O}{\Xi}} is the class-sum lift of such a form, a
function on the adelic group \m{\Gad{\sgO{O}}} of~$\WO{O}$; \m{\JO{O}{\xi}=\pvO{O}{\Har{O}{\xi}}\otimes\Ups{O}} is the
test datum, linear in \m{\xi\in\Vk{n}{\kap{O}}}; and $\Desc{\WO{O}}{\std{E}{n}}{\chO{O}}{\chpO{O}}$ is the theta
descent from the definite to the indefinite group.

\statement{arith}{\phantomsection\label{t:1}\textbf{Theorem 1} (the class-sum lift is a conjugate
descent).\enspace \emph{In Setting~\ref{s:weil}, let $O$ be an auxiliary pair as in Setting~\ref{s:aux} and let
\m{\Xi\in\FormO{O}}. Then, as functions on \m{\Gad{\sgO{O}}},
\begin{align*}
\Th{O}{\Xi}\;=\;{}&\dim\Vk{n}{\kap{O}}\cdot\bigl[\Kz{n}:\KO{O}\bigr]\\
&\cdot\,\conj{\Desc{\WO{O}}{\std{E}{n}}{\chO{O}}{\chpO{O}}\Bigl(\JO{O}{\vplus{\kap{O}}{\paramV{O}{\,\cdot\,}{\,\cdot\,}}},\ \conj{\algF{\paramV{O}{\,\cdot\,}{\,\cdot\,}}{\one{n}}{\Xi}}\Bigr)} ,
\end{align*}
where \m{\paramV{O}{\,\cdot\,}{\,\cdot\,}} is the function \m{(\tau,i)\mapsto\paramV{O}{\tau}{i}}, \m{\vplus{\kap{O}}{\paramV{O}{\,\cdot\,}{\,\cdot\,}}} is the
highest vector for it, \m{\algF{\paramV{O}{\,\cdot\,}{\,\cdot\,}}{\one{n}}{\Xi}} is the scalar function
\m{h\mapsto\gram{\kap{O}}{\vplus{\kap{O}}{\paramV{O}{\,\cdot\,}{\,\cdot\,}}}{\rhoW{\kap{O}}{\eO{O}}(h_{\infty}^{-1})\,\Xi(h_{\mathrm{f}})}} on the
adelic group of~$\std{E}{n}$, and the bars are complex conjugation of functions. The measure of the descent is
normalised as in Setting~\ref{s:weil}.}}
\leanline{\lean{thetaeqsmulstarliftswap}}

\statement{arith}{\phantomsection\label{t:2}\textbf{Theorem 2} (staircases are joint harmonics).\enspace
\emph{Let $O$ be an auxiliary pair as in Setting~\ref{s:aux} and let $l$ be a family of integers $l(\tau,c)$,
indexed by the infinite places~$\tau$ of~$E$ and \m{c\in[r]}, such that \m{\FitsParams{O}{l}}
\textup{(Definition~\ref{d:fits})}. Then there is a family \m{\Phi=(\Phi_{\tau})} of complex polynomials in the
variables \m{X_{ic}}, \m{i\in[n]}, \m{c\in[r]}, one for each infinite place~$\tau$, with these three properties.}
\begin{enumerate}[label=\textup{(\roman*)},leftmargin=2.3em,itemsep=1pt,topsep=2pt]
\item \emph{For every~$\tau$, $\Phi_{\tau}$ is a joint harmonic \textup{(Definition~\ref{d:jh})} for the signs
\m{(\one{n},\sgO{O})}, the weights \m{\bigl(\rmi{O},\,-(\lfloor n/2\rfloor+\wtO{O}{\tau})\bigr)} and the parameters
\m{\bigl(\paramV{O}{\tau}{\,\cdot\,},\,-l(\tau,\,\cdot\,)\bigr)}.}
\item \emph{\m{\Har{O}{\vplus{\kap{O}}{\paramV{O}{\,\cdot\,}{\,\cdot\,}}}=\prod_{\tau}\Phi_{\tau}\bigl((X_{\tau,i,c})_{i,c}\bigr)}: the
polynomial of the highest vector is the product of the $\Phi_{\tau}$, each written in the variables of its place.}
\item \emph{For every~$\tau$ and every \m{c\in[r]}, the lowest $K$-type of the parameter \m{-l(\tau,\,\cdot\,)} for
the signs~$\sgO{O}$ is a staircase weight:}
\[
\lkt{\sgO{O}}{-l(\tau,\,\cdot\,)}(c)=\stw{n}{\sgO{O}}{\operatorname{ord}_{O}}{\kappa^{\pm}_{O}(\tau)}{-(\lfloor n/2\rfloor+\wtO{O}{\tau})}{c},
\]
\emph{where \m{\kappa^{\pm}_{O}(\tau)} is the function with \m{(1,i)\mapsto\kp{O}{\tau}{i}} and
\m{(-1,i)\mapsto\km{O}{\tau}{i}}.}
\end{enumerate}}
\leanline{\lean{existsisJointHarmonic}}

Theorem~1 turns a statement about a finite sum over classes of a definite group into a statement about an
integral of a theta kernel, to which the tools of the theta correspondence apply: the vanishing results of the
third note of this sequence, and the adjunction between \m{\operatorname{Th}_{O}} and the theta lift from
\m{\WO{O}} (§\ref{s:use}). Theorem~2 supplies the archimedean input of that correspondence
(Corollary~\ref{c:edge}): with the datum \m{\JO{O}{\vplus{\kap{O}}{\paramV{O}{\,\cdot\,}{\,\cdot\,}}}}, theta lifts carry forms of the
discrete series type of~$l$ on~$\WO{O}$ to forms of the type of \m{\paramV{O}{\,\cdot\,}{\,\cdot\,}} on the definite group,
and complex conjugates of forms of the type of \m{\paramV{O}{\,\cdot\,}{\,\cdot\,}} back to complex conjugates of
forms of the type of~$l$ (Definition~\ref{d:edge}); accordingly Theorem~1 descends the conjugate
\m{\conj{\algF{\paramV{O}{\,\cdot\,}{\,\cdot\,}}{\one{n}}{\Xi}}}.
The two propositions of §\ref{s:finite} describe the finite part~$\Ups{O}$ of the datum.

\sect{The setting}{s:setting}
\begin{conventions}\phantomsection\label{c:conventions}
Throughout, $E$ is a CM field with a CM type: a number field that is a totally imaginary quadratic extension of a
totally real field, together with the choice, for every infinite place~$\tau$ of~$E$, of a complex embedding that
induces~$\tau$. Complex conjugation of~$E$ and of its completions and adele rings is written \m{z\mapsto\conj{z}};
for a matrix~$X$, \m{X^{*}} is the conjugate transpose and \m{X^{\mathsf{T}}} the transpose. Finite places are
written~$w$; \m{\conj{w}} is the conjugate place, $E_{w}$ the completion, \m{\oo_{w}} its ring of integers and
\m{\FF_{w}} its residue field. The finite adeles of~$E$ are~\m{\AAf}, and \m{\ohat=\prod_{w}\oo_{w}}. The index sets
$[r]$, $[n]$ and~$[r']$ are finite sets with $r$, $n$ and~$r'$ elements; an \m{n\times r} matrix has its rows
indexed by~$[n]$ and its columns by~$[r]$. Where a statement says so, $E$ carries Weil data, described in
Setting~\ref{s:weil}. Signs are elements of \m{\{1,-1\}}, and \m{\one{n}} is the family of signs on~$[n]$ that is
constantly~$1$. The Fock space \m{\mathcal{F}(\varepsilon)} of a family of signs~$\varepsilon$, the vectors
\m{v_{\varepsilon,p}} of polynomials~$p$ and the group $U(\varepsilon)$ are those of the note on Fock vectors.
Letters of the handout's overview that are not displayed with the data as an argument (among them $\Upsilon$,
$R$, $f$, $\varepsilon$, $\beta$, $n$ and~$V$) have, in this note, only the meaning given in its notation list.
\end{conventions}

\begin{setting}[adelic unitary groups, the Weil representation, theta descents; quoted]\label{s:weil}
The following objects are those of the background library; they are recalled, not constructed.

\emph{Weil data.} The field $E$ contains a primitive root of unity of order $5040$ and is equipped with: an
element \m{\jmath\in E} with \m{\conj{\jmath}=-\jmath} whose image under the chosen embedding at every infinite
place has positive imaginary part; a finite idele that normalises the self-duality of~$\ohat$; and a
\emph{splitting character}~$\upsilon$, a unitary Hecke character of~$E$ with
\m{\upsilon(\conj{z})=\upsilon(z)^{-1}}, equal to \m{\prod_{\tau}z_{\tau}/|z_{\tau}|} on the ideles whose finite part
is a unit of~$\ohat$. Its restriction to the ideles of norm one is written~$\upsilon_{1}$.

\emph{Groups.} For a family of signs~$\varepsilon$ on~$[r]$, \m{\Ginf{\varepsilon}} is the product of one copy
of~$U(\varepsilon)$ for each infinite place of~$E$; \m{\Gfin{r}} is the group of matrices
\m{u\in\mathrm{GL}_{r}(\AAf)} with \m{u^{*}u=1}; and \m{\Gad{\varepsilon}=\Ginf{\varepsilon}\times\Gfin{r}}, with
elements \m{u=(u_{\infty},u_{\mathrm{f}})}. The subgroup \m{\Kz{r}=\Gfin{r}\cap\mathrm{GL}_{r}(\ohat)} is open and
compact; a \emph{level} is an open subgroup of it. For a finite place \m{w\ne\conj{w}}, reduction modulo~$w$ is a
homomorphism \m{\red{w}\colon\Kz{r}\to\mathrm{GL}_{r}(\FF_{w})}. A \emph{hermitian space with a point}, of
signs~$\varepsilon$, is a pair of a matrix~$J$ over~$E$ and an element~$\phi$ of \m{\mathrm{GL}_{r}} of the adeles
with \m{\phi^{*}J\phi} equal to \m{\operatorname{diag}(\varepsilon)} at the infinite places and to~$1$ at the finite ones;
its rational points, the \m{\gamma\in\mathrm{GL}_{r}(E)} with \m{\gamma^{*}J\gamma=J}, map to \m{\Gad{\varepsilon}} by
\m{\gamma\mapsto\phi^{-1}\gamma\phi}. The \emph{standard definite space} \m{\std{E}{n}} has \m{J=1} and \m{\phi=1}, and
signs~\m{\one{n}}; its rational points are the unitary matrices over~$E$, and \m{\gamma\mapsto\gamma_{\mathrm{f}}}
denotes their embedding in~\m{\Gfin{n}}. An \emph{automorphic character} is a continuous unitary character of the
adeles of norm one of~$E$ that is trivial on the elements of~$E$ of norm one. It has \emph{weight}
\m{(e_{\tau})_{\tau}}, a family of integers, and \emph{tame part} \m{(\psi_{w})_{w}}, a family of multiplicative
characters of the residue fields~$\FF_{w}$, extended by~$0$ to~$\FF_{w}$ and almost all trivial, if its value at
every~$z$ with integral finite part is \m{\prod_{\tau}z_{\tau}^{e_{\tau}}\prod_{w}\psi_{w}(z\bmod w)}. The
character \m{\upsilon_{1}} has weight~$1$ at every place and trivial tame part.

\emph{The Weil representation.} For hermitian spaces with points $W$, of signs~$\varepsilon$ on~$[r]$, and $V$,
of signs~$\varepsilon'$ on~$[n]$, and automorphic characters $\chi$,~$\chi'$, the library provides a Banach
space \m{\mathcal{S}(W,V)} of \emph{data} on \m{\operatorname{Hom}(W,V)}; a continuous representation
\m{\omega^{W,V}_{\chi,\chi'}} of \m{\Gad{\varepsilon}\times\Gad{\varepsilon'}} on it, in which $(u,h)$ acts by
\m{\chi(\det h)\,\chi'(\det u)^{-1}} times the Weil representation of the unitary group of
\m{\operatorname{Hom}(W,V)} at the map \m{X\mapsto hXu^{-1}}; and a bounded linear functional
\m{\theta_{W,V}} on \m{\mathcal{S}(W,V)}, the \emph{theta functional}, with
\begin{equation}\label{e:rat}
\theta_{W,V}\bigl(\omega^{W,V}_{\chi,\chi'}(\gamma u,\gamma'h)\varphi\bigr)=\theta_{W,V}\bigl(\omega^{W,V}_{\chi,\chi'}(u,h)\varphi\bigr)
\end{equation}
for rational points $\gamma$,~$\gamma'$ of $W$ and~$V$. The \emph{theta descent} of a function~$Y$ on
\m{\Gad{\varepsilon'}} with datum \m{\varphi\in\mathcal{S}(W,V)} is the function on \m{\Gad{\varepsilon}}
\[
\Desc{W}{V}{\chi}{\chi'}(\varphi,Y)(u)=\int_{[V]}\theta_{W,V}\bigl(\omega^{W,V}_{\chi,\chi'}(u,h)\varphi\bigr)\,Y(h)\,\dd h ,
\]
the integral over the quotient of \m{\Gad{\varepsilon'}} by the rational points of~$V$, with the invariant measure
of the library; for \m{V=\std{E}{n}} it is normalised so that \m{\Ginf{\one{n}}\times\Kz{n}} has volume~$1$, the
rational points carrying the counting measure.

\emph{Data.} Let \m{\FockP{\varepsilon}{\varepsilon'}} be the Fock space of the signs
\m{\varepsilon'_{i}\varepsilon_{c}} on the index set of all triples \m{(\tau,i,c)} of an infinite place, an element
of~$[n]$ and an element of~$[r]$, and \m{\pvP{\varepsilon}{\varepsilon'}{p}} the vector in it of a polynomial~$p$
in the variables \m{X_{\tau,i,c}}. For \m{v\in\FockP{\varepsilon}{\varepsilon'}} and \m{f\in\mathcal{S}(W,V)} there is
a datum \m{v\otimes f\in\mathcal{S}(W,V)}, bilinear and bounded in \m{(v,f)}. A \emph{modulus} is a non-zero ideal~$\mathfrak{m}$ of the integers of~$E$ that is
divisible by no pair \m{w,\conj{w}}. For a complex function~$\Psi$ on the \m{n\times r} matrices over the integers
modulo~$\mathfrak{m}$ there is a \emph{finite datum} \m{[\Psi]^{W,V}_{\mathfrak{m}}}: at the places
dividing~$\mathfrak{m}$ it is~$\Psi$ of the reduction of an integral matrix, and elsewhere it is the vacuum, the
characteristic function of the integral matrices. With the Gaussian \m{\varphi_{0}\in\FockP{\varepsilon}{\varepsilon'}}, the datum
\m{\modD{W}{V}{\mathfrak{m}}{\Psi}{\varphi_{0}}} is \m{[\Psi]^{W,V}_{\mathfrak{m}}} itself. Finally
\m{\RPW{\gamma}{\gamma'}}, for families of integers $\gamma$,~$\gamma'$ indexed by the infinite places, is the
representation of \m{\Ginf{\varepsilon}\times\Ginf{\varepsilon'}} on \m{\FockP{\varepsilon}{\varepsilon'}} in which $(u,h)$ acts by
\m{\prod_{\tau}\det(h_{\tau})^{\gamma_{\tau}}\det(u_{\tau})^{-\gamma'_{\tau}}} times the oscillator representation at
\m{X\mapsto hXu^{-1}}.
\end{setting}
\leanline{in the background library: \lnm{CMType}, \lnm{WeilData}, \lnm{WeilData.upsilonU1}, \lnm{GInf},
\lnm{Gf}, \lnm{G}, \lnm{Kf}, \lnm{Kf.reduceAt}, \lnm{Space}, \lnm{stdSpace}, \lnm{stdEmbFin}, \lnm{AutChar},
\lnm{AutChar.HasWeight}, \lnm{pairWeil}, \lnm{pairThetaKer}, \lnm{Space.lift}, \lnm{Space.quotMeasure},
\lnm{Space.integrated},
\lnm{Modulus}, \lnm{vacuumDatum}, \lnm{modDatum}, \lnm{realPairWeil}}

\begin{definition}[level systems; the twist group and the twist lift]\label{d:twist}
A \kw{level system} $\Sigma$ over~$E$ consists of two finite families of finite places of~$E$, the \emph{chosen
places}~$M_{\Sigma}$ and the \emph{twist places}~$Q_{\Sigma}$, such that all these places are distinct and none is
the conjugate of another or of itself. For a level system~$\Sigma$ and an index set~$[n]$ put
\[
\TwistGL{\Sigma}{n}=\prod_{w\in Q_{\Sigma}}\mathrm{GL}_{n}(\FF_{w}),\qquad
\ChGL{\Sigma}{n}=\prod_{w\in M_{\Sigma}}\mathrm{GL}_{n}(\FF_{w}),
\]
finite groups. For \m{A=(A_{w})\in\TwistGL{\Sigma}{n}}, the \kw{twist lift} \m{\twistLift{\Sigma}{A}} is the element
of~\m{\Kz{n}} that the library chooses with \m{\red{w}=1} for \m{w\in M_{\Sigma}} and \m{\red{w}=A_{w}^{-1}} for
\m{w\in Q_{\Sigma}}; its components at the places outside \m{M_{\Sigma}\cup Q_{\Sigma}} and their conjugates
are~$1$.
\end{definition}
\leanline{\lean{TwistGL}\sep \lean{twistLift}\sep in the background library: \lnm{LevelPlaces},
\lnm{LevelPlaces.ChosenGL} (in the unit on the rows), \lnm{liftResidue}}

\begin{setting}[auxiliary pairs; Input from the unit on the rows]\label{s:aux}
Let $E$ carry Weil data (Setting~\ref{s:weil}). An \emph{auxiliary pair}~$O$ over~$E$, for the index sets $[r]$
and~$[n]$, consists of: a level system~$\SigO{O}$, with chosen places~$\MO{O}$ and twist places~$\QO{O}$; a
family of signs~$\sgO{O}$ on~$[r]$ and a hermitian space with a point \m{\WO{O}} of these signs; an injective map
\m{\incO{O}\colon[r]\to[n]}; a function \m{c\mapsto\ordO{O}{c}} from $[r]$ to the natural numbers; natural
numbers \m{\kp{O}{\tau}{i}} and \m{\km{O}{\tau}{i}} for every infinite place~$\tau$ and every natural number~$i$; and
an automorphic character~$\muO{O}$ of weight \m{(\wtO{O}{\tau})_{\tau}} and tame part \m{(\psiO{O}{w})_{w}}, with
\m{\psiO{O}{w}=1} for \m{w\notin\MO{O}\cup\QO{O}}. It is required that \m{r<n}. The following objects are attached
to~$O$.
\begin{enumerate}[label=\textup{(\alph*)},leftmargin=2.2em,itemsep=2pt,topsep=2pt]
\item \m{\rmi{O}} and \m{\rpl{O}} are the numbers of \m{c\in[r]} with \m{\sgO{O}(c)=-1} and with
\m{\sgO{O}(c)=1}. Over any ring, \m{\IncM{O}} is the \m{n\times r} matrix with entry~$1$ at the places
\m{(\incO{O}(c),c)} and $0$ elsewhere; the \emph{$\incO{O}$-rows} of an \m{n\times r} matrix form the \m{r\times r}
matrix of its rows with indices in the image of~$\incO{O}$, and the \emph{$\incO{O}$-block} of an \m{n\times n}
matrix is the \m{r\times r} matrix of its entries with both indices in that image.
\item \emph{Weights.} For an infinite place~$\tau$ and a natural number~$i$,
\begin{align*}
\kap{O}(\tau,i)&=\kp{O}{\tau}{i}+\km{O}{\tau}{n-i}\qquad(\text{with }n-i=0\text{ for }i>n),\\
\eO{O}(\tau)&=-\sum_{1\le i<n}\km{O}{\tau}{i}.
\end{align*}
For a family~$\kappa$ of natural numbers \m{\kappa(\tau,i)}, the space \m{\Vk{n}{\kappa}} is the space of coordinate vectors, in a basis fixed by the library, of the Weyl module of
\m{\prod_{\tau}\mathrm{GL}_{n}(\CC)} given at each place~$\tau$ by the diagram with \m{\kappa(\tau,i)} columns of
length~$i$, \m{0<i<n}; the matrix of a family \m{y=(y_{\tau})} of invertible matrices on it is written
\m{\Sch{\kappa}(y)}. The space has a hermitian
form \m{\gram{\kappa}{\,\cdot\,}{\,\cdot\,}} with an integral Gram matrix, and a complex conjugation
\m{\xi\mapsto\conj{\xi}} of coordinates. The representation \m{\rhoW{\kap{O}}{\eO{O}}} of \m{\Ginf{\one{n}}} on
\m{\Vk{n}{\kap{O}}} lets \m{h=(h_{\tau})} act by
\m{\prod_{\tau}\det(\conj{h_{\tau}})^{\eO{O}(\tau)}\,\Sch{\kap{O}}(\conj{h})}, with \m{\conj{h}} the family of the complex
conjugate matrices; rational points of \m{\std{E}{n}} act through their images at the infinite places.
\item \emph{Level.} \m{\LO{O}} is the group of families \m{(j_{w})} in
\m{\prod_{w\in\MO{O}\cup\QO{O}}\mathrm{GL}_{n}(\FF_{w})} such that \m{j_{w}\IncM{O}=\IncM{O}} for \m{w\in\MO{O}} and,
for \m{w\in\QO{O}}, the entries of~$j_{w}$ in the rows of the image of~$\incO{O}$ and the columns outside it
vanish. Its character~\m{\etaO{O}} is
\[
\etaO{O}(j)=\prod_{w\in\QO{O}}\psiO{O}{w}\bigl(\det(\text{the }\incO{O}\text{-block of }j_{w})\bigr)^{-1}.
\]
The level
\m{\KO{O}} is the group of \m{g'\in\Kz{n}} whose reductions \m{(\red{w}g')_{w}} lie in~\m{\LO{O}}, and
\m{\tetaO{O}(g')=\etaO{O}\bigl((\red{w}g')_{w}\bigr)} is its character.
\item \emph{Forms.} \m{\FormO{O}} is the space of functions \m{\Xi\colon\Gfin{n}\to\Vk{n}{\kap{O}}} with
\[
\Xi(\gamma_{\mathrm{f}}\,h\,g')=\tetaO{O}(g')\,\rhoW{\kap{O}}{\eO{O}}(\gamma)\,\Xi(h)
\]
for all rational points~$\gamma$ of~\m{\std{E}{n}}, all \m{h\in\Gfin{n}} and all \m{g'\in\KO{O}}.
\item \emph{Characters and the Weil representation.} \m{\chO{O}=\upsilon_{1}^{\rmi{O}}} and
\m{\chpO{O}=\muO{O}\,\upsilon_{1}^{\lfloor n/2\rfloor}}. We write \m{\SO{O}}, \m{\omO{O}} and \m{\thO{O}} for
\m{\mathcal{S}(\WO{O},\std{E}{n})}, \m{\omega^{\WO{O},\std{E}{n}}_{\chO{O},\chpO{O}}} and
\m{\theta_{\WO{O},\std{E}{n}}}, and \m{\FockO{O}} and \m{\pvO{O}{p}} for the Fock space
\m{\FockP{\sgO{O}}{\one{n}}} and its vectors \m{\pvP{\sgO{O}}{\one{n}}{p}}.
\item \emph{The finite datum.} For a family \m{X=(X_{w})} of \m{n\times r} matrices over the fields~$\FF_{w}$,
\m{w\in\MO{O}\cup\QO{O}}, put
\[
\PsiO{O}{X}=\prod_{w\in\QO{O}}\psiO{O}{w}\bigl(\det(\text{the }\incO{O}\text{-rows of }X_{w})\bigr)
\]
if \m{X_{w}=\IncM{O}} for all \m{w\in\MO{O}}, and \m{\PsiO{O}{X}=0} otherwise.
With the moduli \m{\lmod{\SigO{O}}{1}=\prod_{w\in\MO{O}}w\prod_{w\in\QO{O}}w} and
\m{\lmod{\SigO{O}}{2}=\prod_{w\in\MO{O}}w^{2}\prod_{w\in\QO{O}}w}, the finite datum of~$O$ is
\[
\Ups{O}=[\Psi_{O}\circ\operatorname{red}]^{\WO{O},\std{E}{n}}_{\lmod{\SigO{O}}{2}},
\]
where \m{\operatorname{red}} sends a matrix modulo \m{\lmod{\SigO{O}}{2}} to the family of its reductions at the
places of \m{\MO{O}\cup\QO{O}}.
\item \emph{The archimedean datum.} Let \m{Y_{\tau,i,a}} be variables indexed by an infinite place, an element
of~$[n]$ and a natural number~$a$. For signs~$s$ on~$[r]$, a function \m{o\colon[r]\to\{0,1,2,\dots\}} and a
sign~$\sigma$, the \emph{letter substitution} \m{\Let{s}{o}{\sigma}} is the homomorphism of algebras from the
polynomials in the \m{Y_{\tau,i,a}} to the polynomials in the \m{X_{\tau,i,c}} that sends \m{Y_{\tau,i,a}} to the sum of
the \m{X_{\tau,i,c}} over the \m{c\in[r]} with \m{s(c)=\sigma} and \m{o(c)=a}. The library provides, for every~$\kappa$,
a linear map \m{\Pol{\kappa}} from \m{\Vk{n}{\kappa}} to the polynomials in the \m{Y_{\tau,i,a}}, and an integral matrix
\m{\mathrm{C}_{O}\colon\Vk{n}{\kap{O}}\to\Vk{n}{\kappa^{+}_{O}}\otimes\Vk{n}{\kappa^{-}_{O}}}, the \emph{mixed Cartan
embedding} \textup{(Q5)}. The linear map \m{\xi\mapsto\Har{O}{\xi}}, from \m{\Vk{n}{\kap{O}}} to the polynomials in the
\m{X_{\tau,i,c}}, is the composite of \m{\mathrm{C}_{O}} with the linear map that sends a tensor \m{\alpha\otimes\beta} to
\m{\Let{\sgO{O}}{\operatorname{ord}_{O}}{1}(\Pol{\kappa^{+}_{O}}\alpha)\cdot\Let{\sgO{O}}{\operatorname{ord}_{O}}{-1}(\Pol{\kappa^{-}_{O}}\beta)}.
The \emph{test datum} is \m{\JO{O}{\xi}=\pvO{O}{\Har{O}{\xi}}\otimes\Ups{O}\in\SO{O}}.
\item \emph{The kernel and the class-sum lift.} For \m{u\in\Gad{\sgO{O}}} and \m{h\in\Gfin{n}}, \m{\cK{O}{u}{h}} is
the linear form on \m{\Vk{n}{\kap{O}}} with
\m{\cK{O}{u}{h}(\xi)=\conj{\thO{O}\bigl(\omO{O}(u,(1,h))\,\JO{O}{\conj{\xi}}\bigr)}}. For \m{\Xi\in\FormO{O}},
\[
\Th{O}{\Xi}(u)=\sum_{[h]}\frac{1}{\#\Gamma_{h}}\;\cK{O}{u}{h}\bigl(\Xi(h)\bigr),
\]
the sum over the double classes of \m{\Gfin{n}} modulo the rational points of \m{\std{E}{n}} on the left and
\m{\KO{O}} on the right, with $\Gamma_{h}$ the finite group of rational points in \m{h\,\KO{O}\,h^{-1}}.
\end{enumerate}
\end{setting}
\leanline{in the unit on the rows: \lnm{AuxPair} with \lnm{AuxPair.κ}, \lnm{AuxPair.detExp}, \lnm{auxLevel},
\lnm{AuxPair.auxLevelChar}, \lnm{AuxPair.levelGroup}, \lnm{AuxPair.levelChar}, \lnm{AuxPair.FormV},
\lnm{AuxPair.a}, \lnm{AuxPair.b}, \lnm{AuxPair.pairThetaKernel}, \lnm{AuxPair.resDetDatum},
\lnm{AuxPair.auxDatum}, \lnm{AuxPair.harmonicPoly}, \lnm{AuxPair.fockDatum}, \lnm{AuxPair.conjKernel},
\lnm{AuxPair.θ}\sep in the background library: \lnm{conjWeylRep}, \lnm{mixedCartan}, \lnm{letterPoly},
\lnm{LevelPlaces.levelModulus}, \lnm{stdClassSum}}

In the terminology of the theta correspondence: \m{(U(\WO{O}),U(\std{E}{n}))} is a dual pair of unitary groups of
ranks \m{r<n} over the totally real subfield, split at the places of the level; \m{\chO{O}} and \m{\chpO{O}} are the
two characters that split the metaplectic cover over the two groups; the datum is a pure tensor whose
archimedean component is \m{\Har{O}{\xi}} times the Gaussian in the Fock model, whose components at the places of
the level are given by~$\Psi_{O}$ and whose other components are unramified. The forms of \m{\FormO{O}} are
algebraic modular forms on a definite unitary group in the sense of Gross.

\begin{definition}[the level indicator]\label{d:indicator}
Let $O$ be an auxiliary pair (Setting~\ref{s:aux}). For an \m{n\times r} matrix~$X$ over the integers of~$E$
modulo \m{\lmod{\SigO{O}}{1}}, with reductions $X_{w}$ at the places \m{w\in\MO{O}\cup\QO{O}}, put
\m{\levelIndicator{O}{X}=1} if \m{X_{w}=\IncM{O}} for every \m{w\in\MO{O}} and the $\incO{O}$-rows of~$X_{w}$ form the
identity matrix for every \m{w\in\QO{O}}; and \m{\levelIndicator{O}{X}=0} otherwise.
\end{definition}
\leanline{\lean{levelIndicator}}

\begin{definition}[the parameter on the definite side]\label{d:paramV}
Let $O$ be an auxiliary pair (Setting~\ref{s:aux}), $\tau$ an infinite place of~$E$ and \m{i\in[n]}. Fix once
and for all a numbering of~$[n]$ by \m{0,\dots,n-1}, and let $\pi(i)$ be the number of~$i$. Then
\[
\paramV{O}{\tau}{i}=2\Bigl(\eO{O}(\tau)+\sum_{\pi(i)<a'<n}\kap{O}(\tau,a')\Bigr)+n-1-2\pi(i).
\]
\end{definition}
\leanline{\lean{paramV}\sep in the background library: \lnm{hcParameter}}

For each~$\tau$, the numbers \m{\paramV{O}{\tau}{i}}, \m{i\in[n]}, are the entries of twice the Harish-Chandra
parameter (highest weight plus half the sum of the positive roots) of the factor at~$\tau$ of the representation
\m{y\mapsto\prod_{\tau}\det(y_{\tau})^{\eO{O}(\tau)}\,\Sch{\kap{O}}(y)} of \m{\Ginf{\one{n}}}, a product of compact unitary
groups of rank~$n$. This representation is the complex conjugate of \m{\rhoW{\kap{O}}{\eO{O}}}, which acts through
the conjugate matrices (Setting~\ref{s:aux}(b)); so the doubled parameter of \m{\rhoW{\kap{O}}{\eO{O}}} itself has
the entries \m{-\paramV{O}{\tau}{i}}. The conjugate is the representation through which the scalar functions
\m{\algF{\paramV{O}{\,\cdot\,}{\,\cdot\,}}{\one{n}}{\Xi}} of Theorem~1 transform under right translation by
\m{\Ginf{\one{n}}}; hence the type \m{\tk{\paramV{O}{\,\cdot\,}{\,\cdot\,}}{\one{n}}} in Corollary~\ref{c:edge}. All
parameters in this note are doubled in this way, which keeps them integral.

\begin{definition}[matched parameters, joint harmonics, staircase weights; quoted]\label{d:jh}
Let $s'$ be a family of signs on~$[n]$, $s$ a family of signs on~$[r]$, and $\gamma$,~$\gamma'$ integers. Consider
the polynomials in the variables \m{X_{ic}}, \m{i\in[n]}, \m{c\in[r]}; the variable \m{X_{ic}} has the sign
\m{s'_{i}s_{c}}. For a variable~$X$ of sign~$1$ let $A_{X}$ be the multiplication by~$X$ and \m{B_{X}=\partial/\partial X};
for a variable of sign~$-1$ let \m{A_{X}=-\partial/\partial X} and let $B_{X}$ be the multiplication by~$X$. The
\emph{row} and \emph{column operators} are, with Kronecker's symbol \m{\delta_{ii'}},
\begin{align*}
R_{ii'}&=\sum_{c\in[r]}A_{X_{ic}}B_{X_{i'c}}+\gamma\,\delta_{ii'}&&(i,i'\in[n]),\\
C_{cc'}&=-\sum_{i\in[n]}B_{X_{ic}}A_{X_{ic'}}+(\gamma'+n)\,\delta_{cc'}&&(c,c'\in[r]).
\end{align*}
They give commuting actions of the Lie algebras of \m{\mathrm{GL}_{n}} and \m{\mathrm{GL}_{r}}: the derivative of the
oscillator representation of the pair, twisted by powers of the determinants.

\emph{Parameters.} A \emph{parameter} for the signs~$s'$ is a family of integers \m{l_{1}(i)}, \m{i\in[n]}; it stands
for twice a Harish-Chandra parameter of the group~$U(s')$. Its \emph{lowest $K$-type} is the family
\begin{align*}
\lkt{s'}{l_{1}}(i)={}&\tfrac12\bigl(l_{1}(i)+\#\{j:l_{1}(j)<l_{1}(i)\}-\#\{j:l_{1}(j)>l_{1}(i)\}\bigr)\\
&-\#\{j:s'_{j}=s'_{i},\,l_{1}(j)<l_{1}(i)\}+\#\{j:s'_{j}=s'_{i},\,l_{1}(j)>l_{1}(i)\},
\end{align*}
the Harish-Chandra parameter plus the half sum of the positive roots minus the sum of the compact positive
roots. The parameter is \emph{spread}, relative to a second index set~$[r]$, if it is injective, all
\m{l_{1}(i)+n-1} are even, and \m{|l_{1}(i)-l_{1}(j)|\ge1000\,(n+r)^{3}} whenever \m{s'_{i}\ne s'_{j}}. The same words
apply to a parameter~$l_{2}$ for the signs~$s$ on~$[r]$, with $n$ and~$r$ exchanged.

\emph{Matching.} Put \m{\xi(i)=l_{1}(i)-2\gamma+r} for \m{i\in[n]} and \m{\xi'(c)=l_{2}(c)-2\gamma'-n} for \m{c\in[r]}.
A pair \m{(i,c)} is \emph{matched} if \m{\xi(i)=-\xi'(c)} and \m{|n-r|<s'_{i}s_{c}\,\xi(i)}. The parameters
\m{(l_{1},l_{2})} \emph{match} for the signs \m{(s',s)} and the weights \m{(\gamma,\gamma')} if both are spread, every
\m{i\in[n]} lies in a matched pair or has \m{|\xi(i)|<n-r}, and every \m{c\in[r]} lies in a matched pair or has
\m{|\xi'(c)|<r-n}.

\emph{Joint harmonics.} A polynomial~$\Phi$ is \emph{jointly highest} of weights \m{(\nu,\nu')} for the parameters
\m{(l_{1},l_{2})} if \m{\Phi\ne0} and
\begin{enumerate}[label=\textup{(\alph*)},leftmargin=2.2em,itemsep=1pt,topsep=2pt]
\item \m{R_{ii}\Phi=\nu(i)\Phi} for all~$i$, and \m{R_{ii'}\Phi=0} whenever \m{s'_{i}=s'_{i'}} and
\m{l_{1}(i')<l_{1}(i)};
\item \m{C_{cc}\Phi=\nu'(c)\Phi} for all~$c$, and \m{C_{cc'}\Phi=0} whenever \m{s_{c}=s_{c'}} and
\m{l_{2}(c')<l_{2}(c)};
\item \m{\partial^{2}\Phi/\partial X_{ic}\partial X_{i'c}=0} for all~$c$ whenever \m{s'_{i}\ne s'_{i'}}, and
\m{\partial^{2}\Phi/\partial X_{ic}\partial X_{ic'}=0} for all~$i$ whenever \m{s_{c}\ne s_{c'}}.
\end{enumerate}
It is a \emph{joint harmonic} for the signs \m{(s',s)}, the weights \m{(\gamma,\gamma')} and the parameters
\m{(l_{1},l_{2})} if these parameters match and $\Phi$ is jointly highest of weights
\m{(\lkt{s'}{l_{1}},\lkt{s}{l_{2}})}.

\emph{Staircase weights.} For signs~$s$ on~$[r]$, a function \m{o\colon[r]\to\{0,1,2,\dots\}}, natural numbers
\m{k(\sigma,i)} for signs~$\sigma$ and natural numbers~$i$, an integer~$\gamma'$ and \m{c\in[r]},
\[
\stw{n}{s}{o}{k}{\gamma'}{c}=\gamma'+n\cdot[s_{c}=-1]-s_{c}\sum_{o(c)<i<n}k(s_{c},i),
\]
where \m{[s_{c}=-1]} is $1$ if \m{s_{c}=-1} and $0$ otherwise.
\end{definition}
\leanline{in the background library: \lnm{FockDualPair.std}, \lnm{oscRows}, \lnm{HCParam.lowestKType},
\lnm{HCParam.HasGap}, \lnm{IsMatch}, \lnm{FockDualPair.IsJointHighest}, \lnm{IsJointHarmonic},
\lnm{staircaseWeight}}

Joint harmonics are the polynomials of lowest degree in which a $K$-type of one group of the pair occurs
together with a $K$-type of the other, in Howe's description of the correspondence in the Fock model; the
library uses them, for parameters that are far apart, to match discrete series representations of the two groups
by their lowest $K$-types. In this note $s'$ is always~\m{\one{n}}, so that conditions on pairs with \m{s'_{i}\ne s'_{i'}}
are empty.

\begin{definition}[parameters that fit an auxiliary pair]\label{d:fits}
Let $O$ be an auxiliary pair (Setting~\ref{s:aux}) and $l$ a family of integers \m{l(\tau,c)}, indexed by the
infinite places~$\tau$ of~$E$ and \m{c\in[r]}. We say that \m{\FitsParams{O}{l}} if the following four conditions
hold.
\begin{enumerate}[label=\textup{(F\arabic*)},leftmargin=2.6em,itemsep=1pt,topsep=2pt]
\item For every~$\tau$, the parameters \m{\bigl(\paramV{O}{\tau}{\,\cdot\,},\,-l(\tau,\,\cdot\,)\bigr)} match for the signs
\m{(\one{n},\sgO{O})} and the weights \m{\bigl(\rmi{O},\,-(\lfloor n/2\rfloor+\wtO{O}{\tau})\bigr)}
(Definition~\ref{d:jh}).
\item For every~$\tau$ and every \m{c\in[r]}, \m{\ordO{O}{c}} is the number of \m{c'\in[r]} with
\m{\sgO{O}(c')=\sgO{O}(c)} and \m{\sgO{O}(c)\,l(\tau,c)<\sgO{O}(c)\,l(\tau,c')}.
\item \m{\kp{O}{\tau}{i}=0} for every~$\tau$ and every \m{i>\rpl{O}}.
\item \m{\km{O}{\tau}{i}=0} for every~$\tau$ and every \m{i>\rmi{O}}.
\end{enumerate}
\end{definition}
\leanline{\lean{FitsParams}\sep in the background library: \lnm{ParamsCorrespond}}

Condition (F2) says that \m{\operatorname{ord}_{O}} numbers the indices of each sign: those of sign~$1$ by
decreasing \m{l(\tau,c)}, those of sign~$-1$ by increasing \m{l(\tau,c)}, and that this numbering is the same at all
infinite places. Conditions (F3) and (F4) bound the lengths of the columns of the diagrams \m{\kappa^{+}_{O}}
and~\m{\kappa^{-}_{O}} by the numbers of positive and negative signs of~$\WO{O}$.

\begin{definition}[forms of a type; edges; quoted]\label{d:edge}
Let $s$ and~$s'$ be families of signs on $[r]$ and~$[n]$. For a function~$\mathbf{e}$ on \m{\Ginf{s}}, a \emph{form
of type~$\mathbf{e}$} on a hermitian space with a point~$W$ of signs~$s$ is a bounded continuous function~$F$ on
\m{\Gad{s}}, left invariant under the rational points of~$W$ and right invariant under an open subgroup of
\m{\Gfin{r}}, with \m{\int_{\Ginf{s}}\mathbf{e}(y)F(uy)\,\dd y=F(u)} for all~$u$ whenever $\mathbf{e}$ is integrable.
For a family~$l$ of parameters \m{l(\tau,\,\cdot\,)} for the signs~$s$, \m{\tk{l}{s}} is the product over the
infinite places of the kernels that the library attaches to the discrete series parameters \m{l(\tau,\,\cdot\,)}:
the formal degree times the complex conjugate of the normalised matrix coefficient. Let $\Omega$ be a map from
\m{\Ginf{s}\times\Ginf{s'}} to the bounded operators of \m{\FockP{s}{s'}}, $\mathbf{e}$ and~$\mathbf{e}'$ functions on
\m{\Ginf{s}} and \m{\Ginf{s'}}, and \m{v\in\FockP{s}{s'}}. Then $v$ is a \emph{half edge} from $\mathbf{e}$
to~$\mathbf{e}'$ for~$\Omega$ if
\begin{enumerate}[label=\textup{(\roman*)},leftmargin=2.3em,itemsep=1pt,topsep=2pt]
\item $\mathbf{e}$ is a cusp kernel of order~$0$ in the sense of the library: an idempotent for convolution,
bounded, integrable against the first power of the norm, and fixed by convolution with some continuous function
of compact support;
\item \m{\int_{\Ginf{s}}\langle v,\Omega(y,1)v\rangle\,\conj{\mathbf{e}(y)}\,\mathbf{e}(1)^{-1}\,\dd y\ne0}, with the inner
product of the Fock space;
\item for every polynomial vector~$v'$ of \m{\FockP{s}{s'}} there is a continuous function~$\alpha$ of compact
support on \m{\Ginf{s'}} such that, for all hermitian spaces with points $W$,~$V$ of signs $s$,~$s'$, all
characters, all data~$f$ and all forms~$Y$ of type~$\mathbf{e}$ on~$W$, if the real points act on the data
\m{v''\otimes f} through~$\Omega$ on~$v''$, then the theta lift~$F$ of~$Y$ to~$V$ with datum \m{v\otimes f} is a form of
type~$\mathbf{e}'$ on~$V$ as soon as it is right invariant under some open subgroup, and the theta lift of~$Y$ with
datum \m{v'\otimes f} is \m{h\mapsto\int_{\Ginf{s'}}\alpha(y)F(hy)\,\dd y}.
\end{enumerate}
It is an \emph{edge} from $\mathbf{e}$ to~$\mathbf{e}'$ for~$\Omega$ if in addition it is a half edge from
\m{\conj{\mathbf{e}'}} to \m{\conj{\mathbf{e}}} for $\Omega$ with its two arguments exchanged.
\end{definition}
\leanline{in the background library: \lnm{Space.IsForm}, \lnm{typeKernel}, \lnm{IsCuspKernel}, \lnm{IsHalfEdge},
\lnm{IsEdge}}

\statement{inksoft}{\textbf{Quoted from the background library.} These facts are taken for granted; references
and their state are in §\ref{s:notes}, note~7. In all of them $W$, $V$, $\varepsilon$, $\varepsilon'$, $\chi$,
$\chi'$ are as in Setting~\ref{s:weil}, and $\omega$ stands for \m{\omega^{W,V}_{\chi,\chi'}}.

\vspace{3pt}\textbf{(Q1)} \emph{Integral elements move the finite datum.} For \m{g\in\Kz{r}}, \m{g'\in\Kz{n}}, a
modulus~$\mathfrak{m}$, a function~$\Psi$ on the \m{n\times r} matrices modulo~$\mathfrak{m}$ and
\m{v\in\FockP{\varepsilon}{\varepsilon'}}:
\m{\omega\bigl((1,g),(1,g')\bigr)\bigl(\modD{W}{V}{\mathfrak{m}}{\Psi}{v}\bigr)=\chi(\det g')\,\chi'(\det g)^{-1}\cdot
\modD{W}{V}{\mathfrak{m}}{\Psi'}{v}}, where \m{\Psi'(X)=\Psi(\conj{g}'^{-1}X\,\conj{g})} and \m{\conj{g}}, \m{\conj{g}'} are the
reductions of $g$,~$g'$ modulo~$\mathfrak{m}$. A function of the reductions at the places dividing~$\mathfrak{m}$
gives the same finite datum for $\mathfrak{m}$ and for every modulus with the same prime divisors.

\vspace{3pt}\textbf{(Q2)} \emph{Characters at integral elements.} If an automorphic character~$\mu$ has a weight
and tame part \m{(\psi_{w})}, trivial outside a finite set~$S$ of places as in Definition~\ref{d:twist}, then
\m{\mu(\det g)=\prod_{w\in S}\psi_{w}(\det\red{w}g)} for every \m{g\in\Kz{r}}. If $\mu$ has weight~$e$ and tame
part~$\psi$, then \m{\mu\,\upsilon_{1}^{a}} has weight \m{(e_{\tau}+a)_{\tau}} and tame part~$\psi$, for every natural
number~$a$.

\vspace{3pt}\textbf{(Q3)} \emph{Real and finite parts act separately.} If $\chi$ has weight~$\gamma$ and
$\chi'$ has weight~$\gamma'$, then for \m{u=(u_{\infty},u_{\mathrm{f}})}, \m{h=(h_{\infty},h_{\mathrm{f}})},
\m{v\in\FockP{\varepsilon}{\varepsilon'}} and a datum~$f$:
\m{\omega(u,h)(v\otimes f)=\bigl(\RPW{\gamma}{\gamma'}(u_{\infty},h_{\infty})v\bigr)\otimes\omega\bigl((1,u_{\mathrm{f}}),(1,h_{\mathrm{f}})\bigr)f}.

\vspace{3pt}\textbf{(Q4)} \emph{The definite group on polynomial vectors.} Let \m{\varepsilon'=\one{n}}, let
$\gamma$ be the constant family with value the number of negative signs of~$\varepsilon$, and $\gamma'$ any
family of integers. For \m{y\in\Ginf{\one{n}}} and a polynomial~$p$ in the \m{X_{\tau,i,c}}:
\m{\RPW{\gamma}{\gamma'}(1,y)\,\pvP{\varepsilon}{\one{n}}{p}=\pvP{\varepsilon}{\one{n}}{y\cdot p}}, where \m{p\mapsto y\cdot p}
is the automorphism of the polynomial algebra that replaces the column \m{(X_{\tau,i,c})_{i}} by
\m{y_{\tau}^{\mathsf{T}}(X_{\tau,i,c})_{i}} if \m{\varepsilon_{c}=1} and by \m{y_{\tau}^{-1}(X_{\tau,i,c})_{i}} if
\m{\varepsilon_{c}=-1}. The map \m{p\mapsto\pvP{\varepsilon}{\one{n}}{p}} is injective.

}
\leanline{in the background library: (Q1) \lnm{pairWeil\_modDatum\_Kf}, \lnm{liftFn\_eq\_residueDatum}\sep (Q2)
\lnm{AutChar.HasWeight.apply\_det\_eq\_prod}, \lnm{AutChar.HasWeight.hasWeight\_rankTwist}\sep (Q3)
\lnm{pairWeil\_integrated\_all}\sep (Q4) \lnm{realPairWeil\_polyVector}, \lnm{polyVector\_injective}}

\statement{inksoft}{\textbf{Quoted from the background library, continued.}

\vspace{3pt}\textbf{(Q5)} \emph{Weyl modules.} For a family \m{y=(y_{\tau})} of invertible complex
\m{n\times n} matrices and every~$\kappa$: \m{\Pol{\kappa}(\Sch{\kappa}(y)\alpha)} is obtained from
\m{\Pol{\kappa}(\alpha)} by replacing each column \m{(Y_{\tau,i,a})_{i}} by \m{y_{\tau}^{\mathsf{T}}(Y_{\tau,i,a})_{i}}. For an
auxiliary pair~$O$ the mixed Cartan embedding satisfies
\[
\mathrm{C}_{O}\,\prod_{\tau}\det(y_{\tau})^{\eO{O}(\tau)}\,\Sch{\kap{O}}(y)=\bigl(\Sch{\kappa^{+}_{O}}(y)\otimes\Sch{\kappa^{-}_{O}}(y^{-\mathsf{T}})\bigr)\,\mathrm{C}_{O},
\]
with \m{y^{-\mathsf{T}}} the family of the inverse transposes, and it sends the vector \m{\vplus{\kap{O}}{l_{1}}} of a
parameter~$l_{1}$ as in (Q9) to a pure tensor of extremal vectors, whose two polynomials are products of minors
of the matrices \m{(Y_{\tau,i,a})}. The matrices \m{\Sch{\kappa}} have polynomial entries with integral
coefficients, so \m{\conj{\Sch{\kappa}(y)\xi}=\Sch{\kappa}(\conj{y})\conj{\xi}}.

\vspace{3pt}\textbf{(Q6)} \emph{Staircases} (Kashiwara--Vergne). Let $s$ be signs on~$[r]$, $l_{1}$ an injective
family of integers on~$[n]$, $l_{2}$ an injective family on~$[r]$, and for \m{c\in[r]} let \m{o(c)} be the number
of~$c'$ with \m{s_{c'}=s_{c}} and \m{s_{c}l_{2}(c')<s_{c}l_{2}(c)}. Let \m{k(\sigma,a)} be natural numbers with
\m{k(\sigma,a)=0} whenever \m{\#\{c:s_{c}=\sigma\}<a<n}, and such that \m{k(1,a)=0} or \m{k(-1,a')=0} whenever
\m{a+a'>n}. For a sign~$\sigma$ and \m{a<n} let \m{D_{\sigma,a}} be the minor of the matrix \m{(X_{ic})} with the
$a$ rows~$i$ for which \m{\sigma\,l_{1}(i)} is largest and the columns~$c$ with \m{s_{c}=\sigma} and \m{o(c)<a}, in
these orders. Then, for all integers $\gamma$,~$\gamma'$, the polynomial
\m{\prod_{\sigma}\prod_{a<n}D_{\sigma,a}^{\,k(\sigma,a)}} is jointly highest for the signs \m{(\one{n},s)}, the weights
\m{(\gamma,\gamma')} and the parameters \m{(l_{1},l_{2})}, of weights \m{(\nu,\nu')} with
\m{\nu'(c)=\stw{n}{s}{o}{k}{\gamma'}{c}} and
\m{\nu(i)=\gamma-\#\{c:s_{c}=-1\}+\sum_{\sigma}\sigma\sum_{a}k(\sigma,a)}, the inner sum over the~$a$ with
\m{\#\{j:\sigma l_{1}(j)>\sigma l_{1}(i)\}<a<n}.

\vspace{3pt}\textbf{(Q7)} \emph{Existence of joint harmonics.} If \m{(l_{1},l_{2})} match for the signs \m{(s',s)}
and the weights \m{(\gamma,\gamma')}, a joint harmonic for them exists.

\vspace{3pt}\textbf{(Q8)} \emph{Uniqueness.} If \m{(l_{1},l_{2})} match, $\Phi_{0}$ is a joint harmonic for them and
$\Phi$ is jointly highest of weights \m{(\lkt{s'}{l_{1}},\nu')} for some~$\nu'$, then $\Phi_{0}$ is a complex
multiple of~$\Phi$.

}
\leanline{in the background library: (Q5) \lnm{PolyRep.MixedCartanEmbedding}, \lnm{mixedCartan},
\lnm{weylRep\_map}\sep (Q6) \lnm{staircase\_harmonic\_highest}\sep (Q7) \lnm{IsMatch.exists\_isJointHarmonic}\sep
(Q8) \lnm{IsMatch.unique}}

\statement{inksoft}{\textbf{Quoted from the background library, continued.}

\vspace{3pt}\textbf{(Q9)} \emph{Parameters of the definite group.} Call a family \m{l_{1}(\tau,i)} a \emph{doubled
Harish-Chandra parameter of} \m{(\kappa,e)} if every \m{l_{1}(\tau,\,\cdot\,)} is injective on~$[n]$ and each of its
values is \m{2(e(\tau)+\sum_{a<a'<n}\kappa(\tau,a'))+n-1-2a} for some \m{a<n}. Then
\m{\lkt{\one{n}}{l_{1}(\tau,\,\cdot\,)}(i)=e(\tau)+\sum_{a<a'<n}\kappa(\tau,a')} with $a$ the number of~$j$ with
\m{l_{1}(\tau,j)>l_{1}(\tau,i)}. The \emph{highest vector} \m{\vplus{\kappa}{l_{1}}\in\Vk{n}{\kappa}} is the vector the
library attaches to the orders of~$[n]$ by decreasing \m{l_{1}(\tau,\,\cdot\,)}; it has real coordinates,
\m{\conj{\vplus{\kappa}{l_{1}}}=\vplus{\kappa}{l_{1}}}, and it is a unit vector:
\m{\gram{\kappa}{\vplus{\kappa}{l_{1}}}{\vplus{\kappa}{l_{1}}}=1}.

\vspace{3pt}\textbf{(Q10)} \emph{From joint harmonics to edges.} Let $[r]$ and~$[n]$ be non-empty, $s$,~$s'$ signs
on them, $\gamma$,~$\gamma'$ families of integers and $l$,~$l'$ families of parameters, all indexed by the
infinite places. If for every~$\tau$ the polynomial~$\Phi_{\tau}$ is a joint harmonic for the signs \m{(s',s)}, the
weights \m{(\gamma_{\tau},-\gamma'_{\tau})} and the parameters \m{(l'(\tau,\,\cdot\,),-l(\tau,\,\cdot\,))}, then the
vector \m{\pvP{s}{s'}{p}} of \m{p=\prod_{\tau}\Phi_{\tau}((X_{\tau,i,c})_{i,c})} is an edge from \m{\tk{l}{s}} to
\m{\tk{l'}{s'}} for \m{\RPW{\gamma}{\gamma'}}.

\vspace{3pt}\textbf{(Q11)} \emph{Automorphy and Schur orthogonality.} (a) Every finite datum is fixed by
\m{\omega((1,g),(1,g'))} for $g$,~$g'$ in suitable levels; if a datum~$\varphi$ is so fixed, then
\m{h\mapsto\theta_{W,V}(\omega(u,h)\varphi)} is continuous, left invariant under the rational points of~$V$ and right
invariant under a level. (b) Let \m{V=\std{E}{n}}, let $l_{1}$ be a doubled Harish-Chandra parameter of
\m{(\kappa,e)}, $\mathbf{K}$ a level in~\m{\Kz{n}}, $\Xi$ an algebraic form of weight \m{(\kappa,e)} as in
Setting~\ref{s:aux}(d), and $\Phi$ a function on \m{\Gad{\one{n}}} with the invariance of~(a) such that
\m{\Phi(y,h)=L_{h}\bigl(\rhoW{\kappa}{e}(y)\vplus{\kappa}{l_{1}}\bigr)} for linear forms~$L_{h}$ on \m{\Vk{n}{\kappa}}.
If \m{h\mapsto L_{h}(\Xi(h))} is right invariant under~$\mathbf{K}$, then
\[
\int_{[\std{E}{n}]}\Phi(h)\,\algF{l_{1}}{\one{n}}{\Xi}(h)\,\dd h
=\frac{1}{\dim\Vk{n}{\kappa}\cdot[\Kz{n}:\mathbf{K}]}\ \sum_{[h]}\frac{1}{\#\Gamma_{h}}\,L_{h}\bigl(\Xi(h)\bigr),
\]
the sum over the double classes modulo the rational points and~$\mathbf{K}$, as in
Setting~\ref{s:aux}\textup{(h)}, and the measure normalised as in Setting~\ref{s:weil}.}
\leanline{in the background library:
(Q9) \lnm{IsHCParam}, \lnm{IsHCParam.lowestKType\_eq}, \lnm{highestVector}, \lnm{star\_highestVector},
\lnm{gramForm\_highestVector\_self}\sep (Q10)
\lnm{isEdge\_of\_isJointHarmonic}\sep (Q11) \lnm{exists\_level\_pairWeil}, \lnm{ThetaKernel.isAutomorphic},
\lnm{Space.integral\_mul\_eq\_classIntegral}, \lnm{classIntegral}, \lnm{classIntegral\_eq\_classSum}; behind (b):
\lnm{Space.integral\_eq\_classIntegral\_of\_isLevel}, \lnm{schur\_orthogonality}}

\sect{The finite datum}{s:finite}
\begin{proposition}[the finite datum is a sum of translates of the level indicator]\label{p:sum}
In Setting~\ref{s:weil}, let $O$ be an auxiliary pair (Setting~\ref{s:aux}). Then
\[
\Ups{O}=\sum_{A\in\TwistGL{\SigO{O}}{r}}\omO{O}\bigl((1,\twistLift{\SigO{O}}{A}),1\bigr)
\Bigl(\modD{\WO{O}}{\std{E}{n}}{\lmod{\SigO{O}}{1}}{\levelIndicator{O}{\,\cdot\,}}{\varphi_{0}}\Bigr),
\]
where \m{(1,\twistLift{\SigO{O}}{A})} is the element of \m{\Gad{\sgO{O}}} with trivial component at the infinite
places.
\end{proposition}
\begin{proof}
By (Q1) both sides are finite data for the modulus \m{\lmod{\SigO{O}}{1}}, the left one of the function
\m{X\mapsto\PsiO{O}{(X_{w})_{w}}}; so it suffices to compare functions of an \m{n\times r} matrix~$X$ modulo
\m{\lmod{\SigO{O}}{1}}. Fix \m{A\in\TwistGL{\SigO{O}}{r}} and put \m{g=\twistLift{\SigO{O}}{A}}; its reduction is~$1$
at the places of~$\MO{O}$ and \m{A_{w}^{-1}} at \m{w\in\QO{O}}. By (Q1) the term of~$A$ is the finite datum of
\m{X\mapsto\chpO{O}(\det g)^{-1}\,\levelIndicator{O}{X\conj{g}}}. By (Q2), applied to \m{\chpO{O}=\muO{O}\upsilon_{1}^{\lfloor n/2\rfloor}}
with the places \m{\MO{O}\cup\QO{O}}, we have
\m{\chpO{O}(\det g)^{-1}=\prod_{w\in\QO{O}}\psiO{O}{w}(\det A_{w})}. And \m{\levelIndicator{O}{X\conj{g}}=1} exactly when
\m{X_{w}=\IncM{O}} for all \m{w\in\MO{O}} and the $\incO{O}$-rows of \m{X_{w}A_{w}^{-1}} form the identity, that is, the
$\incO{O}$-rows of~$X_{w}$ are~$A_{w}$, for all \m{w\in\QO{O}}. So the right side is the finite datum of
\begin{align*}
X\ \longmapsto\ \sum_{A}\ &\prod_{w\in\QO{O}}\psiO{O}{w}(\det A_{w})\\
&\cdot\bigl[X_{w}=\IncM{O}\ (w\in\MO{O}),\ \ \incO{O}\text{-rows of }X_{w}=A_{w}\ (w\in\QO{O})\bigr],
\end{align*}
where the bracket is~$1$ if the conditions in it hold and $0$ otherwise.
If the $\incO{O}$-rows of~$X_{w}$ are invertible for every \m{w\in\QO{O}}, exactly one~$A$ contributes and the
value is \m{\PsiO{O}{(X_{w})_{w}}}. If for some \m{w\in\QO{O}} they are not, no~$A$ contributes, and
\m{\PsiO{O}{(X_{w})_{w}}=0} because \m{\psiO{O}{w}} vanishes at~$0$.
\end{proof}
\leanline{\lean{upsiloneqsum}}

\begin{proposition}[the transformation law of the finite datum]\label{p:law}
In Setting~\ref{s:weil}, let $O$ be an auxiliary pair (Setting~\ref{s:aux}), let \m{g\in\Kz{r}}, \m{g'\in\Kz{n}} and
\m{j\in\LO{O}}, and suppose that \m{\red{w}g'=\operatorname{diag}(\red{w}g,1)\,j_{w}} for every \m{w\in\MO{O}} and
\m{\red{w}g'=j_{w}} for every \m{w\in\QO{O}}, where \m{\operatorname{diag}(g,1)} is the matrix in
\m{\mathrm{GL}_{n}(\FF_{w})} that acts by~$g$ on the basis vectors with indices in the image of~$\incO{O}$ and fixes
the others. Then
\[
\omO{O}\bigl((1,g),(1,g')\bigr)\,\Ups{O}=\Bigl(\prod_{w\in\MO{O}}\psiO{O}{w}(\det\red{w}g)\Bigr)^{-1}\etaO{O}(j)\cdot\Ups{O}.
\]
\end{proposition}
\begin{proof}
By (Q1) the left side is the finite datum of the function
\[
X\ \longmapsto\ \chO{O}(\det g')\,\chpO{O}(\det g)^{-1}\,\PsiO{O}{\bigl((\red{w}g')^{-1}X_{w}\,\red{w}g\bigr)_{w}}.
\]
By (Q2), \m{\chO{O}=\upsilon_{1}^{\rmi{O}}} has trivial tame part, so \m{\chO{O}(\det g')=1}, and
\[
\chpO{O}(\det g)=\prod_{w\in\MO{O}\cup\QO{O}}\psiO{O}{w}(\det\red{w}g).
\]
At \m{w\in\MO{O}}: since $j_{w}$ fixes the basis vectors with indices in the image of~$\incO{O}$,
\[
(\red{w}g')\,\IncM{O}=\operatorname{diag}(\red{w}g,1)\,\IncM{O}=\IncM{O}\,\red{w}g .
\]
Hence \m{(\red{w}g')^{-1}X_{w}\,\red{w}g=\IncM{O}} if and only if \m{X_{w}=\IncM{O}}: the condition at~$\MO{O}$ is
unchanged.

At \m{w\in\QO{O}}: \m{\red{w}g'=j_{w}}, and $j_{w}^{-1}$, like~$j_{w}$, has vanishing entries in the rows of the image
of~$\incO{O}$ and the columns outside it. So the $\incO{O}$-rows of \m{j_{w}^{-1}X_{w}\,\red{w}g} are the
$\incO{O}$-block of~$j_{w}^{-1}$ times the $\incO{O}$-rows of~$X_{w}$ times \m{\red{w}g}, and their determinant is
\m{\det(\incO{O}\text{-block of }j_{w})^{-1}\cdot\det(\incO{O}\text{-rows of }X_{w})\cdot\det\red{w}g}.

Together: the translated function is \m{\etaO{O}(j)\prod_{w\in\QO{O}}\psiO{O}{w}(\det\red{w}g)} times~$\Psi_{O}$, and the
factors at~$\QO{O}$ cancel against those of \m{\chpO{O}(\det g)^{-1}}.
\end{proof}
\leanline{\lean{pairWeilupsilon}\sep in the unit on the rows: \lnm{firstBlockHom}}

For \m{g=1} the hypothesis says that \m{g'\in\KO{O}} with reductions~$j$, and the law reads
\m{\omO{O}(1,(1,g'))\Ups{O}=\tetaO{O}(g')\Ups{O}}: the finite datum is an eigenvector of the level with the character
of the forms. Both families of places contribute inverse character values: at~$\MO{O}$ through \m{\det\red{w}g},
at~$\QO{O}$ through the block of~$j_{w}$, which enters~\m{\etaO{O}} with an inverse.

\sect{The archimedean datum}{s:arch}
\begin{lemma}[the parameter of the definite side]\label{l:hc}
Let $O$ be an auxiliary pair (Setting~\ref{s:aux}). Then \m{(\tau,i)\mapsto\paramV{O}{\tau}{i}} is a doubled
Harish-Chandra parameter of \m{(\kap{O},\eO{O})} in the sense of \textup{(Q9)}.
\end{lemma}
\begin{proof}
By Definition~\ref{d:paramV}, the value at \m{(\tau,i)} is the number of (Q9) for \m{a=\pi(i)}, and
\m{a\mapsto2(\eO{O}(\tau)+\sum_{a<a'<n}\kap{O}(\tau,a'))+n-1-2a} is strictly decreasing, hence injective.
\end{proof}
\leanline{\lean{isHCParamparamV}\sep in the background library: \lnm{isHCParam\_hcParameter}}

\begin{lemma}[letter substitutions]\label{l:letters}
Let $s$ be a family of signs on~$[r]$, let \m{o\colon[r]\to\{0,1,2,\dots\}}, let $\gamma'$ be a family of integers
indexed by the infinite places, \m{y\in\Ginf{\one{n}}}, $\sigma$ a sign and $F$ a polynomial in the variables
\m{Y_{\tau,i,a}}. Let $\gamma$ be the constant family whose value is the number of \m{c\in[r]} with \m{s(c)=-1}. Then
\[
\RPW{\gamma}{\gamma'}(1,y)\ \pvP{s}{\one{n}}{\Let{s}{o}{\sigma}(F)}=\pvP{s}{\one{n}}{\Let{s}{o}{\sigma}(F^{y,\sigma})},
\]
where \m{F^{y,\sigma}} is obtained from~$F$ by replacing each column \m{(Y_{\tau,i,a})_{i}} by
\m{y_{\tau}^{\mathsf{T}}(Y_{\tau,i,a})_{i}} if \m{\sigma=1}, and by \m{y_{\tau}^{-1}(Y_{\tau,i,a})_{i}} if \m{\sigma=-1}.
\end{lemma}
\begin{proof}
By (Q4) the left side is the vector of \m{y\cdot\Let{s}{o}{\sigma}(F)}. Both \m{F\mapsto y\cdot\Let{s}{o}{\sigma}(F)} and
\m{F\mapsto\Let{s}{o}{\sigma}(F^{y,\sigma})} are homomorphisms of algebras, so it suffices to compare them on a
variable \m{Y_{\tau,i,a}}. Let $z_{\tau}$ be \m{y_{\tau}^{\mathsf{T}}} or \m{y_{\tau}^{-1}} according to~$\sigma$. The
columns~$c$ that occur in \m{\Let{s}{o}{\sigma}(Y_{\tau,i,a})} have \m{s(c)=\sigma}, so the first map gives
\m{\sum_{c}\sum_{i'}z_{\tau}(i,i')X_{\tau,i',c}}, the sums over \m{i'\in[n]} and over the~$c$ with \m{s(c)=\sigma} and
\m{o(c)=a}; the second gives the same double sum in the other order.
\end{proof}
\leanline{\lean{realPairWeilfockLetters}}

\begin{lemma}[the definite group acts on the harmonic polynomials through the weight]\label{l:har}
Let $O$ be an auxiliary pair (Setting~\ref{s:aux}), $\gamma'$ a family of integers indexed by the infinite
places, \m{y\in\Ginf{\one{n}}} and \m{\xi\in\Vk{n}{\kap{O}}}. Let $\gamma$ be the constant family with value~\m{\rmi{O}}. Then
\[
\RPW{\gamma}{\gamma'}(1,y)\ \pvO{O}{\Har{O}{\conj{\xi}}}=\pvO{O}{\Har{O}{\conj{\rhoW{\kap{O}}{\eO{O}}(y)\,\xi}}}.
\]
\end{lemma}
\begin{proof}
By Setting~\ref{s:aux}(b) and (Q5),
\[
\conj{\rhoW{\kap{O}}{\eO{O}}(y)\,\xi}=\prod_{\tau}\det(y_{\tau})^{\eO{O}(\tau)}\,\Sch{\kap{O}}(y)\,\conj{\xi},
\]
and the mixed Cartan embedding sends this vector to
\m{\bigl(\Sch{\kappa^{+}_{O}}(y)\otimes\Sch{\kappa^{-}_{O}}(y^{-\mathsf{T}})\bigr)\mathrm{C}_{O}\conj{\xi}}. Write
\m{\mathrm{C}_{O}\conj{\xi}} as a sum of pure tensors \m{\alpha\otimes\beta}. By (Q5), in the notation of
Lemma~\ref{l:letters},
\[
\Pol{\kappa^{+}_{O}}\bigl(\Sch{\kappa^{+}_{O}}(y)\alpha\bigr)=(\Pol{\kappa^{+}_{O}}\alpha)^{y,1},\qquad
\Pol{\kappa^{-}_{O}}\bigl(\Sch{\kappa^{-}_{O}}(y^{-\mathsf{T}})\beta\bigr)=(\Pol{\kappa^{-}_{O}}\beta)^{y,-1}.
\]
So \m{\Har{O}{\conj{\rhoW{\kap{O}}{\eO{O}}(y)\xi}}} is the sum of the products
\[
\Let{\sgO{O}}{\operatorname{ord}_{O}}{1}\bigl((\Pol{\kappa^{+}_{O}}\alpha)^{y,1}\bigr)\cdot\Let{\sgO{O}}{\operatorname{ord}_{O}}{-1}\bigl((\Pol{\kappa^{-}_{O}}\beta)^{y,-1}\bigr).
\]
By Lemma~\ref{l:letters}, (Q4) and the injectivity in (Q4), each factor is \m{y\cdot{}} the corresponding factor of
\m{\Har{O}{\conj{\xi}}}; as \m{p\mapsto y\cdot p} is a homomorphism of algebras, the sum is
\m{y\cdot\Har{O}{\conj{\xi}}}. By (Q4) once more, its vector is the left side.
\end{proof}
\leanline{\lean{realPairWeilharmonicPoly}}

\begin{proof}[Proof of Theorem 2]
Fix an infinite place~$\tau$. Put \m{l_{1}=\paramV{O}{\tau}{\,\cdot\,}} and \m{l_{2}=-l(\tau,\,\cdot\,)}, and let
\m{k(1,i)=\kp{O}{\tau}{i}}, \m{k(-1,i)=\km{O}{\tau}{i}}, \m{\gamma=\rmi{O}}, \m{\gamma'=-(\lfloor n/2\rfloor+\wtO{O}{\tau})}.

\emph{Step 1: the staircase.} The hypotheses of (Q6) hold for \m{s=\sgO{O}}: $l_{1}$ is injective by
Lemma~\ref{l:hc}, and $l_{2}$ is injective because it is spread, by (F1). The function~$o$ of (Q6) is
\m{\operatorname{ord}_{O}}, by (F2): the inequality \m{s_{c}l_{2}(c')<s_{c}l_{2}(c)} is \m{s_{c}\,l(\tau,c)<s_{c}\,l(\tau,c')}.
The vanishing of \m{k(\sigma,a)} for \m{\#\{c:s_{c}=\sigma\}<a} is (F3) and (F4). If \m{a+a'>n}, then
\m{a+a'>r=\rpl{O}+\rmi{O}} because \m{r<n}, so \m{a>\rpl{O}} or \m{a'>\rmi{O}}, and \m{k(1,a)=0} or \m{k(-1,a')=0}. Let $\Phi_{\tau}$ be the
polynomial of (Q6). It is jointly highest of weights \m{(\nu,\nu')} with \m{\nu'(c)=\stw{n}{\sgO{O}}{\operatorname{ord}_{O}}{k}{\gamma'}{c}}.

\emph{Step 2: its weight on the definite side is the lowest $K$-type.} Let \m{i\in[n]} and let $a$ be the
number of~$j$ with \m{l_{1}(j)>l_{1}(i)}; then \m{n-1-a} is the number of~$j$ with \m{-l_{1}(j)>-l_{1}(i)}. As
\m{\gamma=\rmi{O}}, (Q6) gives
\m{\nu(i)=\sum_{a<a'<n}\kp{O}{\tau}{a'}-\sum_{n-1-a<a'<n}\km{O}{\tau}{a'}}. On the other hand, by Lemma~\ref{l:hc} and
(Q9),
\begin{align*}
\lkt{\one{n}}{l_{1}}(i)&=\eO{O}(\tau)+\sum_{a<a'<n}\kap{O}(\tau,a')\\
&=-\sum_{1\le a'<n}\km{O}{\tau}{a'}+\sum_{a<a'<n}\kp{O}{\tau}{a'}+\sum_{1\le a'<n-a}\km{O}{\tau}{a'},
\end{align*}
which is the same number. So \m{\nu=\lkt{\one{n}}{l_{1}}}.

\emph{Step 3: its weight on the indefinite side.} By (F1) the parameters \m{(l_{1},l_{2})} match, so by (Q7)
there is a joint harmonic~$\Phi_{0}$ for them, and by (Q8) and Step~2, $\Phi_{0}$ is a complex multiple
of~$\Phi_{\tau}$. For \m{c\in[r]} the operator \m{C_{cc}} multiplies $\Phi_{0}$ by \m{\lkt{\sgO{O}}{l_{2}}(c)} and every
multiple of~$\Phi_{\tau}$ by \m{\nu'(c)}; since \m{\Phi_{0}\ne0}, the two numbers agree. This is~(iii). It also shows that
$\Phi_{\tau}$ is jointly highest of weights \m{(\lkt{\one{n}}{l_{1}},\lkt{\sgO{O}}{l_{2}})}, hence a joint harmonic:
this is~(i).

\emph{Step 4: the polynomial of the highest vector.} By (Q5) the mixed Cartan embedding sends
\m{\vplus{\kap{O}}{\paramV{O}{\,\cdot\,}{\,\cdot\,}}} to a pure tensor of extremal vectors, and their polynomials are products of minors of
the matrices \m{(Y_{\tau,i,a})}: for the diagram \m{\kappa^{+}_{O}}, the product over~$\tau$ and \m{a<n} of the
\m{\kp{O}{\tau}{a}}-th power of the minor with the first~$a$ columns and the $a$ rows of largest \m{l_{1}}; for
\m{\kappa^{-}_{O}}, the same with \m{\km{O}{\tau}{a}} and the $a$ rows of smallest~$l_{1}$. The letter substitutions
\m{\Let{\sgO{O}}{\operatorname{ord}_{O}}{\pm1}} replace the column of index~$a'$ by the column \m{X_{\tau,\,\cdot\,,c}}
with \m{\sgO{O}(c)=\pm1} and \m{\ordO{O}{c}=a'}, which is unique when it exists, by (F2). This turns the two products into
\m{\prod_{\tau}\prod_{a}D_{1,a}^{\,k(1,a)}} and \m{\prod_{\tau}\prod_{a}D_{-1,a}^{\,k(-1,a)}} in the variables of each place, and
their product is \m{\prod_{\tau}\Phi_{\tau}}. This is~(ii).
\end{proof}
\leanline{\lean{existsisJointHarmonic}}

\begin{corollary}[the vector of the highest polynomial is an edge]\label{c:edge}
In Setting~\ref{s:weil}, let $O$ be an auxiliary pair (Setting~\ref{s:aux}) with \m{[r]} not empty, and let $l$ be
a family of integers \m{l(\tau,c)} with \m{\FitsParams{O}{l}}. Let $\gamma$ be the constant family with
value~\m{\rmi{O}} and \m{\gamma'(\tau)=\lfloor n/2\rfloor+\wtO{O}{\tau}}. Then the vector
\m{\pvO{O}{\Har{O}{\vplus{\kap{O}}{\paramV{O}{\,\cdot\,}{\,\cdot\,}}}}} is an edge \textup{(Definition~\ref{d:edge})} from \m{\tk{l}{\sgO{O}}} to
\m{\tk{\paramV{O}{\,\cdot\,}{\,\cdot\,}}{\one{n}}} for the representation \m{\RPW{\gamma}{\gamma'}}, where \m{\paramV{O}{\,\cdot\,}{\,\cdot\,}} is the family
\m{(\tau,i)\mapsto\paramV{O}{\tau}{i}}.
\end{corollary}
\begin{proof}
Since \m{[r]} is not empty and \m{\incO{O}} maps it into~$[n]$, the set~$[n]$ is not empty. Theorem~2(i) and~(ii)
give the hypothesis of (Q10) for \m{s'=\one{n}}, \m{s=\sgO{O}} and \m{l'=\paramV{O}{\,\cdot\,}{\,\cdot\,}}, and its conclusion is the
assertion.
\end{proof}
\leanline{\lean{isEdge}}

The weights of \m{\RPW{\gamma}{\gamma'}} in the corollary are those of the characters \m{\chO{O}} and \m{\chpO{O}}, by
(Q2). So, by (Q3), the real points act on the data \m{v\otimes f} of the Weil representation \m{\omO{O}} through
\m{\RPW{\gamma}{\gamma'}}, and the corollary applies to the theta lifts between \m{\WO{O}} and \m{\std{E}{n}} with the
test datum.

\sect{The class-sum lift as a descent}{s:descent}
\begin{lemma}[the transformation law of the test datum]\label{l:datum}
In Setting~\ref{s:weil}, let $O$ be an auxiliary pair (Setting~\ref{s:aux}), \m{y\in\Ginf{\one{n}}}, \m{g'\in\KO{O}}
and \m{\xi\in\Vk{n}{\kap{O}}}. Then
\[
\omO{O}\bigl(1,(y,g')\bigr)\,\JO{O}{\conj{\xi}}=\tetaO{O}(g')\cdot\JO{O}{\conj{\rhoW{\kap{O}}{\eO{O}}(y)\,\xi}}.
\]
\end{lemma}
\begin{proof}
By (Q2), the weight of \m{\chO{O}} is \m{\rmi{O}} at every place, and the weight of \m{\chpO{O}} at~$\tau$ is
\m{\lfloor n/2\rfloor+\wtO{O}{\tau}}. By (Q3) the left side is the tensor product of \m{\RPW{\gamma}{\gamma'}(1,y)\pvO{O}{\Har{O}{\conj{\xi}}}}, with
these weights, and of \m{\omO{O}(1,(1,g'))\Ups{O}}. The first is \m{\pvO{O}{\Har{O}{\conj{\rhoW{\kap{O}}{\eO{O}}(y)\xi}}}} by
Lemma~\ref{l:har}; the second is \m{\tetaO{O}(g')\Ups{O}} by Proposition~\ref{p:law} with \m{g=1} and~$j$ the family
of the reductions of~$g'$.
\end{proof}
\leanline{\lean{pairWeilfockDatum}}

\begin{lemma}[the kernel is equivariant]\label{l:kernel}
In Setting~\ref{s:weil}, let $O$ be an auxiliary pair (Setting~\ref{s:aux}), $\gamma$ a rational point of
\m{\std{E}{n}}, \m{g'\in\KO{O}}, \m{v\in\Vk{n}{\kap{O}}}, \m{u\in\Gad{\sgO{O}}} and \m{h\in\Gfin{n}}. Then
\[
\cK{O}{u}{\gamma_{\mathrm{f}}\,h\,g'}\bigl(\tetaO{O}(g')\,\rhoW{\kap{O}}{\eO{O}}(\gamma)\,v\bigr)=\cK{O}{u}{h}(v).
\]
\end{lemma}
\begin{proof}
Put \m{\xi=\rhoW{\kap{O}}{\eO{O}}(\gamma)v}. As \m{\xi\mapsto\JO{O}{\conj{\xi}}} is conjugate-linear, the left side is
\m{\tetaO{O}(g')} times the complex conjugate of
\m{\thO{O}\bigl(\omO{O}(u,(1,\gamma_{\mathrm{f}}hg'))\JO{O}{\conj{\xi}}\bigr)}. By Lemma~\ref{l:datum} with \m{y=1},
\[
\omO{O}\bigl(1,(1,g')\bigr)\,\JO{O}{\conj{\xi}}=\tetaO{O}(g')\,\JO{O}{\conj{\xi}};
\]
the values of \m{\tetaO{O}} are roots of unity, so the two factors \m{\tetaO{O}(g')} and \m{\conj{\tetaO{O}(g')}}
cancel. It remains to show that
\[
\thO{O}\bigl(\omO{O}(u,(1,\gamma_{\mathrm{f}}h))\,\JO{O}{\conj{\xi}}\bigr)=\thO{O}\bigl(\omO{O}(u,(1,h))\,\JO{O}{\conj{v}}\bigr).
\]
The rational point~$\gamma$ is \m{(\gamma_{\infty},\gamma_{\mathrm{f}})} in \m{\Gad{\one{n}}}, and
\m{(1,\gamma_{\mathrm{f}}h)=\gamma\cdot(\gamma_{\infty}^{-1},h)}. By \eqref{e:rat} the left side is
\[
\thO{O}\Bigl(\omO{O}\bigl(u,(1,h)\bigr)\,\omO{O}\bigl(1,(\gamma_{\infty}^{-1},1)\bigr)\,\JO{O}{\conj{\xi}}\Bigr),
\]
and by Lemma~\ref{l:datum} with \m{g'=1} the inner datum is
\m{\JO{O}{\conj{\rhoW{\kap{O}}{\eO{O}}(\gamma_{\infty}^{-1})\xi}}=\JO{O}{\conj{v}}}.
\end{proof}
\leanline{\lean{conjKernelequivariant}}

\begin{proof}[Proof of Theorem 1]
Write
\[
v^{+}=\vplus{\kap{O}}{\paramV{O}{\,\cdot\,}{\,\cdot\,}},\qquad F_{\Xi}=\algF{\paramV{O}{\,\cdot\,}{\,\cdot\,}}{\one{n}}{\Xi}.
\]
Fix \m{u\in\Gad{\sgO{O}}} and put
\m{\Phi(h)=\conj{\thO{O}\bigl(\omO{O}(u,h)\JO{O}{v^{+}}\bigr)}} for \m{h\in\Gad{\one{n}}}. The complex conjugate of the
descent at~$u$ is \m{\int_{[\std{E}{n}]}\Phi(h)F_{\Xi}(h)\,\dd h}: the datum is conjugated with the value of the theta
functional, and the two conjugations of~$F_{\Xi}$ cancel. We apply (Q11)(b) with \m{\mathbf{K}=\KO{O}},
\m{\kappa=\kap{O}}, \m{e=\eO{O}}, \m{l_{1}=\paramV{O}{\,\cdot\,}{\,\cdot\,}} (Lemma~\ref{l:hc}) and \m{L_{h}=\cK{O}{u}{h}}.

The function~$\Phi$ has the invariance of (Q11)(a), because \m{\JO{O}{v^{+}}} is a polynomial vector tensor the
finite datum~\m{\Ups{O}}. For \m{h=(y,h_{\mathrm{f}})}, Lemma~\ref{l:datum} with \m{g'=1}, and \m{v^{+}=\conj{v^{+}}} from
(Q9), give
\[
\Phi(y,h_{\mathrm{f}})=\conj{\thO{O}\bigl(\omO{O}(u,(1,h_{\mathrm{f}}))\,\JO{O}{\conj{\rhoW{\kap{O}}{\eO{O}}(y)v^{+}}}\bigr)}
=\cK{O}{u}{h_{\mathrm{f}}}\bigl(\rhoW{\kap{O}}{\eO{O}}(y)v^{+}\bigr).
\]
The function \m{h\mapsto\cK{O}{u}{h}(\Xi(h))} on \m{\Gfin{n}} is right invariant under \m{\KO{O}}: by
Setting~\ref{s:aux}(d), \m{\Xi(hg')=\tetaO{O}(g')\Xi(h)}, and Lemma~\ref{l:kernel} with \m{\gamma=1} gives
\[\cK{O}{u}{hg'}\bigl(\tetaO{O}(g')\,\Xi(h)\bigr)=\cK{O}{u}{h}\bigl(\Xi(h)\bigr).\]
So (Q11)(b) applies:
\[
\int_{[\std{E}{n}]}\Phi(h)\,F_{\Xi}(h)\,\dd h=\frac{1}{\dim\Vk{n}{\kap{O}}\cdot[\Kz{n}:\KO{O}]}\sum_{[h]}\frac{1}{\#\Gamma_{h}}\,\cK{O}{u}{h}\bigl(\Xi(h)\bigr),
\]
and the sum is \m{\Th{O}{\Xi}(u)} by Setting~\ref{s:aux}(h).
\end{proof}
\leanline{\lean{thetaeqsmulstarliftswap}}

Lemma~\ref{l:kernel} with general~$\gamma$ shows that the summand of the class sum does not depend on the
representative of the double class, so that \m{\Th{O}{\Xi}} is well defined by Setting~\ref{s:aux}(h); the proof of
Theorem~1 uses the lemma only for \m{\gamma=1}.

\sect{Where the results are used}{s:use}
\emph{In the unit on the rows.} The auxiliary pairs of the proof are the pairs \m{(W_{j},V_{j})} of the $j$-th
theta lift of the chain, \m{1\le j\le m}, and a parameter that fits them is provided by the unit on the exchange
identity. For a form~$\Xi$ on the definite group with \m{\Th{O}{\Xi}=0}, Theorem~1 shows that the theta descent of
\m{\conj{\algF{\paramV{O}{\,\cdot\,}{\,\cdot\,}}{\one{n}}{\Xi}}} to~\m{\WO{O}} with the datum
\m{\JO{O}{\vplus{\kap{O}}{\paramV{O}{\,\cdot\,}{\,\cdot\,}}}} vanishes. Corollary~\ref{c:edge} extends this to the data with any
polynomial vector at the infinite places. The third note of this sequence then extends it to all data that are
unramified outside the level, and to the spaces \m{\WO{O}} plus hyperbolic planes; this is how forms that are not
lifts are shown to have vanishing theta periods.

\emph{In the unit on the exchange identity.} If \m{\FitsParams{O}{l}} and $\WO{O}$ is not zero, the map
\m{\operatorname{Th}_{O}} is adjoint, for the Petersson products, to the theta lift from \m{\WO{O}} to the definite
group with the datum \m{\JO{O}{\vplus{\kap{O}}{\paramV{O}{\,\cdot\,}{\,\cdot\,}}}}: Theorem~1 and Corollary~\ref{c:edge} are the two inputs of this
adjunction, and Lemma~\ref{l:datum} gives the level and the weight of the lifted form.
Proposition~\ref{p:sum} is used where the datum~$\Ups{O}$ is compared with data of full level.

\emph{In the unit on the twists.} The image of \m{\operatorname{Th}_{O}} is identified with a space of theta lifts by
the same two results.

\sect{Notes on the formal text}{s:notes}
\begin{enumerate}[label=\textup{\arabic*.},leftmargin=1.6em,itemsep=5pt]
\item \emph{What is a definition and what is a theorem.} The class sum and the kernel of Setting~\ref{s:aux}(h)
are the definition of \m{\operatorname{Th}_{O}} in the unit on the rows; the integral with its scalar is the theorem
of this unit. In Theorem~1 the basic descent is bilinear in the value of the theta functional and in its input
function; the input is the complex conjugate of~$F_{\Xi}$ and the output is conjugated again. The formal statement
writes the scalar as the number of elements of the index set of the basis of \m{\Vk{n}{\kap{O}}} times the index
\m{[\Kz{n}:\KO{O}\cap\Kz{n}]}; since \m{\KO{O}\subseteq\Kz{n}} this is the index of Theorem~1.
\item \emph{Arguments read from the text, not printed by Lean.} The file declares its variables twice: the
field, the index sets and the pair are introduced at the top and again before the definition of
\m{\paramV{O}{\tau}{i}}, with the letters of the index sets changed. The program that reads the arguments of a
declaration from the text cannot tell the two declarations of~$E$ apart. For the eight lemmas after the second
declaration the dossier gives the list of arguments that Lean's rule for variables predicts, by hand, and marks
it as not printed by Lean, since nothing was built. I read the two variable lines and the eight statements, and
the lists agree with what I find: each statement takes the field with its two structures, the two index sets,
the pair where it occurs, the Weil data from the lemma on edges onwards, and its own arguments, in this order.
The statements of this note account for exactly these arguments. A print by Lean of the eight signatures is
still to be made.
\item \emph{Index sets.} In the formal text $[r]$ and~$[n]$ are arbitrary finite types, $r$ and~$n$ their numbers
of elements, and the numbering~$\pi$ of Definition~\ref{d:paramV} is a bijection chosen by the library. In the
first part of the file the two index sets are called by other letters; the note uses $[r]$ for the indefinite
and $[n]$ for the definite side throughout. The twist group of Definition~\ref{d:twist} is defined for any index
set; it is used for~$[r]$.
\item \emph{The modulus.} The finite datum is defined with the modulus \m{\lmod{\SigO{O}}{2}}, in which the chosen
places occur squared, and Proposition~\ref{p:sum} is stated with \m{\lmod{\SigO{O}}{1}}. The function~$\Psi_{O}$
depends only on the reductions at the places, so the two data agree; the formal proof begins with this change
of modulus.
\item \emph{Steps carried out by automation.} The reader of the unit marked the central steps of
Propositions \ref{p:sum} and~\ref{p:law} as descriptions of what the simplifier computes, not replayed goal by
goal. I read these lines. The lemmas they name are those used in the proofs above: the action of integral
elements on finite data (Q1), the values of the characters (Q2), the reductions of the twist lift, the two
rules for products with matrices in the level (rows of a product with a matrix of block triangular shape, and
\m{\operatorname{diag}(g,1)\IncM{O}=\IncM{O}g}), and the case distinction on the invertibility of the tested blocks,
which is explicit in the Lean text. The proofs above are complete as mathematics; that the simplifier's steps
are these was not replayed. In Lemma~\ref{l:har} the formal proof uses the compatibility of \m{\Pol{\kappa}} with the
matrices \m{\Sch{\kappa}} through rules of the library that are not named in the text of the proof; the form in
which (Q5) states it, with a transpose, is the one that makes Lemmas \ref{l:letters} and~\ref{l:har} agree, and I
did not read its proof in the library. Likewise Step~4 of the proof of Theorem~2 describes a computation that
the formal proof leaves to the simplifier, after unfolding the highest vector, the letter substitutions and the
polynomials of extremal vectors.
\item \emph{A weaker hypothesis would do, and a stronger conclusion holds.} Theorem~2 uses of (F1) only that the
parameters match; of the requirement \m{r<n} only that \m{i+i'>n} forces \m{i>\rpl{O}} or \m{i'>\rmi{O}}. The formal
proof of Theorem~2 does not use Weil data, and the theorem is stated without them; Corollary~\ref{c:edge} has them
because edges are defined with theta lifts.
\item \emph{References for the quoted facts.} The catalogue of the handout and the second reader's comparison
with the literature give the following; where no place is named, the \emph{reference is to be confirmed}.
(Q6): M.~Kashiwara, M.~Vergne, \emph{On the Segal--Shale--Weil representations and harmonic polynomials},
Invent.~Math.~44 (1978), Chapter~III, Proposition~6.1 with formula~(6.2) and Theorem~6.3, and M.~G.~Davidson,
\emph{The harmonic representation of U(p,q) and its connection with the generalized unit disk}, Pacific
J.~Math.~129 (1987), §2, Theorem~2.9; the formal statement is a \emph{special case} (the explicit products of
minors under the condition on the sizes, in the orders given by the parameters, with the shifts of the weights
by the vacuum and the determinant characters). (Q7): R.~Howe, \emph{Transcending classical invariant theory},
J.~Amer.~Math.~Soc.~2 (1989), §3, Lemmas 3.2 and~3.3, and A.~Paul and P.~Trapa, \emph{One-dimensional
representations of U(p,q) and the Howe correspondence}, §3, Lemma~3.2(1)--(2), with formulas (3.3), (3.4)
and~(3.6); a \emph{special case} (existence only for matched parameters with large gaps, in the Fock model): the
correspondence of the compact types and its formula for the minimal degree give the joint harmonic (the second
reader's verdict). (Q8): Howe, loc.~cit.; J.-S.~Li, Duke Math.~J.~61 (1990); A.~Paul, J.~Funct.~Anal.~159 (1998); a
\emph{special case} (uniqueness up to a scalar only for matched parameters with large gaps, in the Fock model);
places not confirmed. (Q10): A.~Ichino, \emph{Theta lifting for discrete series
representations of real unitary groups}, §7.2, Proposition~7.2, and §7.4, Lemma~7.8; Howe, loc.~cit., Theorem~1;
the relation is \emph{unresolved}: the published results give the correspondence of the compact types, and the
two-sided analytic conclusions of an edge need further argument; the catalogue describes the formal statement
as a \emph{weaker form} of the correspondence of discrete series, without Howe duality and without
non-vanishing of the lift. The Weil representation of the pair and its theta functional, (Q1) to (Q4): the
library's file names A.~Weil, \emph{Sur certains groupes d'opérateurs unitaires}, which the artifact's list of
references marks ``work found, no place given''; reference to be confirmed. (Q5), (Q9) and (Q11): no reference
recorded; (Q11)(b) is Schur's orthogonality for the compact group \m{\Ginf{\one{n}}} together with the definition of
the measure: its factor is the product of \m{1/\dim\Vk{n}{\kappa}}, from the orthogonality relation for the unit
vector \m{\vplus{\kappa}{l_{1}}} on a group of volume~$1$, and of \m{1/[\Kz{n}:\mathbf{K}]}, the volume
of~$\mathbf{K}$ when \m{\Kz{n}} has volume~$1$. The two normalisations are those of the library: the header of its
lemma behind (Q11)(b) says ``vol~$\mathbf{K}_{0}=1$'', and (Q9) has the unit vector.
\end{enumerate}

\subhead{How this note was written.} From the dossier of the unit: sixteen members and seven structured proofs.
The program check of the unit's dossier covered \UnitSteps{} leaf steps of \UnitProofs{} proofs, found the
statements at their lines and compared the quotations of the imported declarations with the source. I read the
whole Lean file of the unit and compared every statement with the statements above, argument by argument. In the
unit on the rows I read the definition of an auxiliary pair and of all objects of Setting~\ref{s:aux}. In the
background library I read the definitions behind Setting~\ref{s:weil} and Definitions \ref{d:jh} and~\ref{d:edge}
(groups, spaces with a point, characters with weight and tame part, the Weil representation of a pair and the
theta functional as they are packaged, the lift, finite data, the row and column operators, lowest $K$-types,
gaps, matching, joint harmonics, staircase weights, cusp kernels, forms of a type, half edges and edges) and the
statements of the lemmas behind (Q1) to (Q11); I read none of their proofs, and I did not read the construction of
the space of data, of the theta functional, of the quotient measure or of the Weyl modules. The normalisation of
the measure in Setting~\ref{s:weil} is read from the statements and headers of the library's lemmas behind
(Q11)(b) and of its Schur orthogonality, which holds as stated only for total mass~$1$ of \m{\Ginf{\one{n}}}. The second reader of the artifact read the
unit independently; her remarks on the two conjugations of Theorem~1, on the inverse in the character~\m{\etaO{O}}
and on the signed ordering in (F2) agree with the text above. The proofs follow the formal ones; those of the two
propositions of §\ref{s:finite} and of Lemma~\ref{l:kernel} are written out from the definitions and are longer than
the structured proofs, and the remarks after the statements and §\ref{s:use} are mine. For §\ref{s:use} I read the
statements of the consuming lemmas.

\head{Appendix A. Concordance with the Lean artifact}
{\fontsize{8.6}{11.6}\selectfont\renewcommand{\leanfont}{\ttfamily\fontsize{7.7}{10}\selectfont\color{leancol}}
\setlength{\LTleft}{0pt}\setlength{\LTright}{\fill}\setlength{\tabcolsep}{0pt}
\begin{longtable}{@{}p{30mm}p{56mm}p{43mm}r@{}}
\textbf{In this note} & \textbf{Lean declaration} & \textbf{File} & \textbf{Lines}\\\midrule\endhead
\ConcordanceRows
\end{longtable}}

\subhead{The two theorems as they are written in the artifact.}
\quoted{thetaeqsmulstarliftswap}
\quoted{existsisJointHarmonic}
\subhead{Legend.} In the formal text the auxiliary pair is {\leanfont O}; {\leanfont O.signs}, {\leanfont O.W},
{\leanfont O.incl}, {\leanfont O.ord}, {\leanfont O.colsPlus}, {\leanfont O.colsMinus}, {\leanfont O.μ},
{\leanfont O.weight}, {\leanfont O.ψ} are \m{\sgO{O}}, \m{\WO{O}}, \m{\incO{O}}, \m{\operatorname{ord}_{O}},
\m{\kappa^{+}_{O}}, \m{\kappa^{-}_{O}}, \m{\muO{O}}, \m{\operatorname{wt}_{O}}, \m{\psi_{O}}; {\leanfont O.κ} and
{\leanfont O.detExp} are \m{\kap{O}} and \m{\eO{O}}; {\leanfont O.a} and {\leanfont O.b} are \m{\chO{O}} and \m{\chpO{O}};
{\leanfont negCount O.signs} is \m{\rmi{O}}; {\leanfont O.auxDatum} is \m{\Ups{O}}, {\leanfont O.harmonicPoly} is
\m{\operatorname{Har}_{O}}, {\leanfont O.fockDatum} is \m{\operatorname{J}_{O}}, {\leanfont O.conjKernel} is
\m{\mathcal{C}_{O}} and {\leanfont O.θ} is \m{\operatorname{Th}_{O}}; {\leanfont O.levelGroup}, {\leanfont O.levelChar}
and {\leanfont O.auxLevelChar} are \m{\KO{O}}, \m{\tetaO{O}} and \m{\etaO{O}}; {\leanfont Kf} is \m{\mathbf{K}_{0}}; the
letters {\leanfont k}, {\leanfont k'} for its elements are $g$,~$g'$ here, the letter {\leanfont g} for a point of
the adelic group of~$\WO{O}$ is~$u$, the sign {\leanfont c} of the lemma on letters is~$\sigma$, the signs {\leanfont t} of the library's
definitions on~$[n]$ are~$s'$ (the letter \emph{t} is kept in this write-up for the step of the ranks), and the
infinite places, written {\leanfont w} in part of the file, are~$\tau$. {\leanfont Space.lift} with the swapped
kernel is the descent, {\leanfont star} on functions is complex conjugation, {\leanfont ofAlgForm} gives
\m{\algF{l_{1}}{\one{n}}{\Xi}}, and {\leanfont realPairWeil γ γ'} is \m{\RPW{\gamma}{\gamma'}}.

\head{Notation}
\begingroup\fontsize{8.6}{10.9}\selectfont\raggedright
\setlength{\columnsep}{9mm}\setlength{\parskip}{0pt}
\begin{multicols}{2}
\nt{$E$, $\tau$, $w$, $\conj{w}$}{a CM field with a CM type; an infinite place; a finite place and its conjugate}
\nt{$E_{w}$, $\oo_{w}$, $\FF_{w}$, $\AAf$, $\ohat$}{completion, its integers, residue field; finite adeles; their
maximal compact subring}
\nt{$[r]$, $[n]$, $[r']$, $r$, $n$}{finite index sets; their numbers of elements (ranks of the indefinite and the
definite space)}
\nt{$i$, $i'$, $j$; $c$, $c'$; $a$, $a'$}{indices in $[n]$ ($j$ is also an element of \m{\LO{O}}); indices in
$[r]$; natural numbers (column indices of diagrams, sizes of minors, exponents)}
\nt{$X^{*}$, $X^{\mathsf{T}}$, $\conj{z}$}{conjugate transpose, transpose, complex conjugate}
\nt{$\jmath$, $\upsilon$, $\upsilon_{1}$}{the Weil data: a purely imaginary element, the splitting character, its
restriction to norm one}
\nt{$\varepsilon$, $\varepsilon'$, $s$, $s'$, $\sigma$, $\one{n}$}{families of signs on index sets; a sign; the
constant family~$1$ on $[n]$}
\nt{$U(\varepsilon)$, $\Ginf{\varepsilon}$, $\Gfin{r}$, $\Gad{\varepsilon}$}{the unitary group of the signs;
its product over the infinite places; the finite adelic unitary group; their product}
\nt{$u$, $h$, $y$}{points of the adelic group of the indefinite, resp.\ definite, space, with components
$u_{\infty}$, $u_{\mathrm{f}}$, $h_{\infty}$, $h_{\mathrm{f}}$; $y$ a point of a real group, mostly of \m{\Ginf{\one{n}}} (in (Q5), $y$ is any
family of invertible complex matrices)}
\nt{$\Kz{n}$, $g$, $g'$, $\red{w}$, $\mathbf{K}$}{the integral points of \m{\Gfin{n}}; elements of such groups;
reduction modulo~$w$; a level}
\nt{$W$, $V$, $J$, $\phi$, $\std{E}{n}$}{hermitian spaces with a point; the matrix of the form and the point; the standard
definite space}
\nt{$\gamma$, $\gamma_{\infty}$, $\gamma_{\mathrm{f}}$}{a rational point and its components (in
§\ref{s:descent}); otherwise $\gamma$, $\gamma'$ are weights: integers or families of integers}
\nt{$\chi$, $\chi'$, $\mu$, $e$, $\psi_{w}$}{automorphic characters; a weight; a tame part}
\nt{$\mathcal{S}(W,V)$, $\omega^{W,V}_{\chi,\chi'}$, $\theta_{W,V}$}{data, the Weil representation of the pair, the
theta functional}
\nt{$\varphi$, $f$, $v\otimes f$, $\varphi_{0}$}{data; the datum with archimedean part~$v$; the Gaussian}
\nt{$\Desc{W}{V}{\chi}{\chi'}(\varphi,Y)$, $Y$, $\dd h$}{the theta descent of a function~$Y$ from $V$ to~$W$; the
invariant measure}
\nt{$\FockP{\varepsilon}{\varepsilon'}$, $\pvP{\varepsilon}{\varepsilon'}{p}$, $p$}{the Fock space of the pair; the
vector of a polynomial~$p$}
\nt{$X_{\tau,i,c}$, $X_{ic}$, $Y_{\tau,i,a}$, $X$}{the variables of the Fock model; those of one place; the matrix
variables of Weyl modules; $X$ also a matrix over a residue ring, with reductions $X_{w}$}
\nt{$\mathfrak{m}$, $[\Psi]^{W,V}_{\mathfrak{m}}$, $\Psi$, $\Psi'$}{a modulus; the finite datum of a function on
matrices modulo~$\mathfrak{m}$}
\nt{$\RPW{\gamma}{\gamma'}$, $\Omega$}{the representation of the real points on the Fock space with weights
$\gamma$,~$\gamma'$; a map to operators of the Fock space}
\nt{$\Sigma$, $M_{\Sigma}$, $Q_{\Sigma}$, $S$}{a level system, its chosen and twist places; a finite set of places}
\nt{$\TwistGL{\Sigma}{n}$, $\ChGL{\Sigma}{n}$, $A$}{the products of \m{\mathrm{GL}_{n}(\FF_{w})} over the twist, resp.\
chosen, places; an element of the first}
\nt{$\twistLift{\Sigma}{A}$}{the twist lift of~$A$}
\nt{$O$}{an auxiliary pair; $\SigO{O}$, $\MO{O}$, $\QO{O}$ its level system and places}
\nt{$\WO{O}$, $\sgO{O}$, $\rpl{O}$, $\rmi{O}$}{its indefinite space, signs, numbers of positive and negative signs}
\nt{$\incO{O}$, $\IncM{O}$, $\ordO{O}{c}$}{its embedding of index sets, the matrix of the embedding, its numbering}
\nt{$\kp{O}{\tau}{i}$, $\km{O}{\tau}{i}$, $\kap{O}$, $\eO{O}$}{column numbers; the weight on the definite side}
\nt{$\muO{O}$, $\wtO{O}{\tau}$, $\psiO{O}{w}$}{its character, with weight and tame part}
\nt{$\Vk{n}{\kappa}$, $\Sch{\kappa}$, $\gram{\kappa}{\xi}{\eta}$, $\Pol{\kappa}$}{coordinate vectors of the Weyl
module, its matrices, its hermitian form, its polynomial realisation}
\nt{$\rhoW{\kappa}{e}$, $\xi$, $\eta$, $\alpha$, $\beta$, $v$}{the weight representation; vectors of Weyl modules
($\alpha$ also a function of compact support, $v$ also a Fock vector)}
\nt{$\LO{O}$, $\etaO{O}$, $\KO{O}$, $\tetaO{O}$}{the level in the residue groups, its character; the level,
its character}
\nt{$\FormO{O}$, $\Xi$}{algebraic forms on the definite group; one of them}
\nt{$\chO{O}$, $\chpO{O}$, $\SO{O}$, $\omO{O}$, $\thO{O}$, $\FockO{O}$, $\pvO{O}{p}$}{the characters, data, Weil
representation, theta functional, Fock space and polynomial vectors of~$O$}
\nt{$\PsiO{O}{X}$, $\lmod{\Sigma}{1}$, $\lmod{\Sigma}{2}$, $\Ups{O}$}{the function of residue matrices, the two
moduli, the finite datum}
\nt{$\levelIndicator{O}{X}$}{the level indicator}
\nt{$\Let{s}{o}{\sigma}$, $o$, $\mathrm{C}_{O}$, $F$}{the letter substitution; a numbering; the mixed Cartan
embedding; a polynomial in the \m{Y_{\tau,i,a}} (in Definition~\ref{d:edge} and §\ref{s:descent} also a form)}
\nt{$\Har{O}{\xi}$, $\JO{O}{\xi}$}{the harmonic polynomial and the test datum of~$\xi$}
\nt{$\cK{O}{u}{h}$, $\Th{O}{\Xi}$, $\Gamma_{h}$}{the kernel, the class-sum lift, the stabiliser of a class}
\nt{$\paramV{O}{\tau}{i}$, $\paramV{O}{\,\cdot\,}{\,\cdot\,}$, $\pi$}{the parameter on the definite side; the numbering of $[n]$}
\nt{$l$, $l_{1}$, $l_{2}$, $l'$}{parameters (doubled Harish-Chandra parameters)}
\nt{$\FitsParams{O}{l}$}{Definition~\ref{d:fits}}
\nt{$A_{X}$, $B_{X}$, $R_{ii'}$, $C_{cc'}$, $\delta_{ii'}$}{raising and lowering operators; row and column
operators; Kronecker's symbol}
\nt{$\lkt{s'}{l_{1}}$, $\xi(i)$, $\xi'(c)$}{the lowest $K$-type; the centred parameters}
\nt{$\Phi$, $\Phi_{\tau}$, $\Phi_{0}$, $\nu$, $\nu'$}{polynomials (joint harmonics); weights ($\Phi$ is also an
automorphic function in (Q11))}
\nt{$\stw{n}{s}{o}{k}{\gamma'}{c}$, $k$, $D_{\sigma,a}$}{the staircase weight; column numbers; the minors of (Q6)}
\nt{$\tk{l}{s}$, $\mathbf{e}$, $\mathbf{e}'$}{the discrete series kernel of~$l$; functions on the real groups}
\nt{$\vplus{\kappa}{l_{1}}$, $v^{+}$, $\algF{l_{1}}{\one{n}}{\Xi}$, $F_{\Xi}$, $L_{h}$}{the highest vector; the scalar
function of a form; linear forms}
\nt{$F^{y,\sigma}$, $y\cdot p$, $z_{\tau}$}{substitutions in polynomials; a matrix}
\nt{$z$}{an idele or adele (Setting~\ref{s:weil})}
\end{multicols}
\endgroup

\end{document}
